\documentclass[14pt, a4paper]{extarticle}

% --------------------------------------
%            ПОДКЛЮЧЕНИЕ ПАКЕТОВ
% --------------------------------------

\usepackage[
a4paper,
left=3cm,
right=1cm,
top=2.5cm,
bottom=2cm,
%includefoot
]{geometry}
\usepackage[russian]{babel}
\usepackage{xurl}
\usepackage{float}
\usepackage{placeins}
\usepackage{fontspec}
\usepackage[justification=centering]{caption}
\usepackage{listings}
\usepackage{multicol}
\usepackage{xcolor}
\usepackage{enumitem}
\usepackage{titlesec}
\usepackage{tocloft}
\usepackage{indentfirst}
\usepackage{setspace}
\usepackage{listings}
\usepackage{etoolbox}
\usepackage{graphicx}
\usepackage{titleps}
\usepackage{lipsum}
\usepackage[export]{adjustbox}
\usepackage{newtxmath}
\usepackage{array}
\usepackage{csquotes}
\usepackage{longtable}
\usepackage{ltablex}
\usepackage{pdfpages}
\usepackage{ragged2e}
\usepackage{pgf}
\usepackage{microtype}
\usepackage{totcount}
%\usepackage{showframe}
\usepackage{needspace}
\usepackage{afterpage}
\usepackage{tikz}
\usepackage[bottom]{footmisc}
\usetikzlibrary{tikzmark}
\usepackage[hypertexnames=false]{hyperref}
\usepackage[
  backend=biber,
  style=numeric,
  sorting=customsort,
  defernumbers=true
]{biblatex}
% ----------------------------------------
\keepXColumns
% --------------------------------------
%            НАСТРОЙКИ ЛИСТИНГОВ
% --------------------------------------

\renewcommand{\lstlistingname}{Листинг}

\newcommand{\vaultpath}{\py{__import__("os").environ.get("VAULT_PATH") or "../vault_diploma"}}

%\newcommand{\listingsgeometry}{}

% --------------------------------------
%            СТАРЫЕ ЦВЕТА 
% --------------------------------------

\definecolor{codegray}{rgb}{0.97,0.97,0.97}
\definecolor{lightgray}{rgb}{0.98,0.98,0.98}
\definecolor{white}{rgb}{1,1,1}
\definecolor{codebg}{rgb}{0.95,0.95,0.95}
\definecolor{keywordcolor}{rgb}{0.0,0.0,0.7}
\definecolor{stringcolor}{rgb}{0.6,0.0,0.0}
\definecolor{commentcolor}{rgb}{0.0,0.5,0.0}

% --------------------------------------

% --------------------------------------
%            НОВЫЕ ЦВЕТА 
% --------------------------------------

\definecolor{bg}{RGB}{250,250,250}
\definecolor{kw}{RGB}{0,92,197}
\definecolor{com}{RGB}{0,128,0}
\definecolor{str}{RGB}{163,21,21}
\definecolor{text}{RGB}{25,25,25}

% --------------------------------------

\lstdefinelanguage{bashcustom}{
    keywords={cd, ls, pwd, echo, export, python},
    sensitive=true,
    morecomment=[l]{\#},
    morestring=[b]",
    morestring=[b]',
}

\newfontfamily\listingmono[
  Path=./fonts/latin-modern-mono/,
  Extension=.otf,
  UprightFont=lmmono10-regular,
  BoldFont=lmmonolt10-bold,
  ItalicFont=lmmonolt10-oblique,
  BoldItalicFont=lmmonolt10-boldoblique,
  RawFeature={fallback=monocyrfallback},
  BoldFeatures={RawFeature={fallback=monocyrfallback}},
  ItalicFeatures={RawFeature={fallback=monocyrfallback}},
  BoldItalicFeatures={RawFeature={fallback=monocyrfallback}},
]{lmmono10-regular}

\lstset{
    showstringspaces=false,
    language=Python,
    basicstyle=\listingmono\small,
    numberstyle=\small,
    stepnumber=1,
    frame=none,
    breaklines=true,
    breakatwhitespace=false,
    tabsize=4,
    captionpos=b,
    backgroundcolor=\color{codegray},
    columns=fixed,
    keepspaces=true,
    xrightmargin=5pt,
    extendedchars=true,
	keywordstyle=\bfseries,    
    numbers=left,
    xleftmargin=1.1em,
    numbersep=0.5em,
    aboveskip=\baselineskip,
    belowskip=\baselineskip,
    escapeinside={'(*@}{@*)'},
    inputpath=\vaultpath
}

% --------------------------------------
%            НАСТРОЙКИ МЕТОК 
% --------------------------------------

\newcommand{\codecallout}[1]{
  \hspace{-2.5em}
  \tikz[baseline=(n.base)]{
    \node[
      circle,
      fill=black!85,
      text=white,
      inner sep=1.2pt,
      minimum size=1.25em,
      font=\scriptsize\bfseries
    ] (n) {#1};
  }
  \hspace{0.2em}
}

\newcommand{\codecalloutref}[1]{\ref{#1}}

\newcommand{\codecalloutlistingref}[1]{\ref{#1@listing}}

\makeatletter
\newcommand{\codecalloutlabel}[1]{
  \refstepcounter{codecallout}
  \label{#1}
  \protected@write\@auxout{}{
    \string\newlabel{#1@listing}{{\thelstlisting}{\thepage}{}{}{}}
  }%
  \codecallout{\thecodecallout}
}
\makeatother

\newcommand{\calloutfullref}[1]{
  \codecalloutref{#1}
  (см. с.~\pageref{#1}, листинг~\codecalloutlistingref{#1}, приложение~\listingsappendix{})
}

\newcommand{\calloutfullreffootnote}[1]{\footnote{Метка \codecalloutref{#1}, листинг \ref{#1@listing}, приложение~\listingsappendix{}, с.~\pageref{#1}}}

\newcommand{\calloutfullreffootnotetextonly}[1]{
   \codecalloutref{#1}, листинг~\codecalloutlistingref{#1}, приложение~\listingsappendix{}, с.~\pageref{#1}
}

% --------------------------------------

\lstdefinestyle{modern}{
    backgroundcolor=\color{bg},
    basicstyle=\ttfamily\scriptsize\color{text},
    keywordstyle=\color{kw}\bfseries,
    commentstyle=\color{com}\itshape,
    stringstyle=\color{str},
    numbers=left,
    numberstyle=\scriptsize\texttt,
    frame=single,
    rulecolor=\color{gray!30},
    breaklines=true,
    showstringspaces=false,
    tabsize=2,
    numbersep=1.2em,
    xleftmargin=1.1em,
    xrightmargin=0.2cm,
}

\lstdefinestyle{modernNoNumbers}{
    backgroundcolor=\color{bg},
    basicstyle=\ttfamily\scriptsize,
    keywordstyle=\color{kw}\bfseries,
    commentstyle=\color{com}\itshape,
    stringstyle=\color{str},
    numbers=none,
    frame=single,
    rulecolor=\color{gray!30},
    breaklines=true,
    showstringspaces=false,
    tabsize=2,
    numbersep=0.5em,
    xleftmargin=1.1em,
    xrightmargin=0.2cm,
}


\lstdefinestyle{bashcommand}{
    backgroundcolor=\color{bg},
    basicstyle=\listingmono\scriptsize\color{text},
    keywordstyle=\bfseries,
    commentstyle=\color{com}\itshape,
    stringstyle=\color{str},
    numbers=none,
    frame=single,
    rulecolor=\color{gray!30},
    breaklines=true,
    showstringspaces=false,
    tabsize=2,
    numbersep=0.5em,
    xrightmargin=0.3cm,
    framexleftmargin=0.5em,
    framextopmargin=0.5em,
    framexbottommargin=0.5em,
    framexrightmargin=0.5em,
}

% --------------------------------------
%            НАСТРОЙКИ СОДЕРЖАНИЯ
% --------------------------------------

\addto\captionsrussian{\renewcommand{\contentsname}{}}

% Заголовок содержания
\renewcommand{\cfttoctitlefont}{\hfill\normalsize}
\renewcommand{\cftaftertoctitle}{\hfill}
\setlength{\cftbeforetoctitleskip}{0pt}
\setlength{\cftaftertoctitleskip}{1.5\baselineskip}

% Отступы уровней
\setlength{\cftsecindent}{0pt}
\setlength{\cftsubsecindent}{0.8cm}
\setlength{\cftsubsubsecindent}{1.9cm}

% Место под номер раздела
\setlength{\cftsecnumwidth}{0.7cm}
\setlength{\cftsubsecnumwidth}{1cm}
\setlength{\cftsubsubsecnumwidth}{1.5cm}

% Вертикальные интервалы
\setlength{\cftbeforesecskip}{0.45\baselineskip}
\setlength{\cftbeforesubsecskip}{0.25\baselineskip}
\setlength{\cftbeforesubsubsecskip}{0.15\baselineskip}

% Шрифт
\renewcommand{\cftsecfont}{\normalsize}
\renewcommand{\cftsubsecfont}{\normalsize}
\renewcommand{\cftsubsubsecfont}{\normalsize}
\renewcommand{\cftsecpagefont}{\normalsize}
\renewcommand{\cftsubsecpagefont}{\normalsize}
\renewcommand{\cftsubsubsecpagefont}{\normalsize}

% Без точек-лидеров
\renewcommand{\cftdot}{}
\renewcommand{\cftsecleader}{\hfill}
\renewcommand{\cftsubsecleader}{\hfill}
\renewcommand{\cftsubsubsecleader}{\hfill}

\renewcommand{\UrlFont}{\rmfamily}
% --------------------------------------
%       НАСТРОЙКИ ОТСТУПОВ ОТ ПУНКТОВ
% --------------------------------------

\titleformat{\section}
{\normalsize}
{\thesection}
{0.5em}
{}

\titleformat{\subsection}
{\normalsize\bfseries}
{\thesubsection}
{0.5em}
{}

\titleformat{\subsubsection}
{\normalsize}
{\thesubsubsection}
{0.5em}
{}

\titlespacing*{\section}
{\parindent}
{0pt}
{0pt}

\titlespacing*{\subsection}
{\parindent}
{0pt}
{0pt}

\titlespacing*{\subsubsection}
{\parindent}
{0pt}
{0pt}


% ------------------------------------------------
% МАКРОС ДЛЯ GENERIC ГЛАВЫ (ВВЕДЕНИЕ, ЗАКЛЮЧЕНИЕ)
% ------------------------------------------------

\newcommand{\genericsection}[1]{
\phantomsection
\addcontentsline{toc}{section}{#1}


\centerline{\MakeUppercase{#1}}

%\vspace*{-0.2cm}
}
% --------------------------------------
%            НАСТРОЙКИ СЧЕТЧИКОВ
% --------------------------------------

%  ------------- РЕГИСТРАЦИЯ ---------
\regtotcounter{table}
\regtotcounter{lstlisting}
\newcounter{bibcount}
\newcounter{appendixcount}
\newcounter{codecallout}

\newcommand{\newappendix}{
  \stepcounter{appendixcount}
}

%  ------------- СЧЕТЧИК ПРИЛОЖЕНИЙ ---------
\makeatletter
\AtEndDocument{%
  \immediate\write\@auxout{%
    \string\gdef\string\totalappendices{\arabic{appendixcount}}%
  }%
}
\makeatother

%  ------------- СЧЕТЧИКИ ДО БИБЛИОГРАФИИ ---------
\makeatletter
\newcommand{\recordPreBibliographyCounters}{%
  \immediate\write\@auxout{%
    \string\gdef\string\tablesBeforeBib{\number\value{table}}%
  }%
  \immediate\write\@auxout{%
    \string\gdef\string\listingsBeforeBib{\number\value{lstlisting}}%
  }%
}
\makeatother

%  ------------- СЧЕТЧИК БИБЛИОГРАФИИ ---------
\AtBeginBibliography{%
  \setcounter{bibcount}{0}%
  \recordPreBibliographyCounters%
}

\makeatletter
\newcommand{\recordBibliographyTotals}{%
  \immediate\write\@auxout{%
    \string\gdef\string\totalbibitems{\number\value{bibcount}}%
  }%
}
\makeatother

\AtEveryBibitem{%
  \stepcounter{bibcount}%
  \recordBibliographyTotals%
}

\providecommand{\totalbibitems}{0}
\providecommand{\tablesBeforeBib}{0}
\providecommand{\listingsBeforeBib}{0}
\providecommand{\totalappendices}{0}

% -----------------------------------

% --------------------------------------
%  НАСТРОЙКИ МНОЖЕСТВЕННЫХ ФОРМ ДЛЯ РЕФЕРАТА
% --------------------------------------

\directlua{
  russian_plural = russian_plural or {}

  function russian_plural.form(number, one, few, many)
    local n = math.abs(tonumber(number) or 0)
    local last_two = n % 100
    local last = n % 10

    if last_two >= 11 and last_two <= 14 then
      return many
    end

    if last == 1 then
      return one
    end

    if last >= 2 and last <= 4 then
      return few
    end

    return many
  end
}

\newcommand{\russianPlural}[4]{%
  \directlua{
    tex.sprint(russian_plural.form(
      "\luaescapestring{#1}",
      "\luaescapestring{#2}",
      "\luaescapestring{#3}",
      "\luaescapestring{#4}"
    ))
  }%
}

\newcommand{\russianTables}[1]{%
  \russianPlural{#1}{таблицу}{таблицы}{таблиц}%
}

\newcommand{\russianAppendices}[1]{%
  \russianPlural{#1}{приложение}{приложения}{приложений}%
}

\newcommand{\russianSources}[1]{%
  \russianPlural{#1}{источник}{источника}{источников}%
}

\newcommand{\russianListings}[1]{%
  \russianPlural{#1}{листинг}{листинга}{листингов}%
}

% --------------------------------------
%          НАСТРОЙКИ БИБЛИОГРАФИИ
% --------------------------------------

\addbibresource{bibliography.bib}

\setcounter{biburlnumpenalty}{9000}
\setcounter{biburlucpenalty}{9000}
\setcounter{biburllcpenalty}{9000}

\def\UrlBreaks{\do\/\do-\do\_\do\.\do\?\do\&\do\=\do\#}

\DeclareSortingTemplate{customsort}{
  \sort{\field{sortkey}}
}

\DeclareBibliographyDriver{misc}{
  \printfield{note}
  \finentry
}

\DeclareFieldFormat{labelnumberwidth}{#1}

\defbibenvironment{bibliography}
  {\list
     {\printtext[labelnumberwidth]{\printfield{labelnumber}}}
     {
       \setlength{\labelwidth}{2em}
       \setlength{\labelsep}{0.5em}
       \setlength{\leftmargin}{0cm}
       \setlength{\itemindent}{2.2cm}
       \setlength{\listparindent}{0pt}
       \setlength{\itemsep}{0pt}
       \setlength{\parsep}{0pt}
     }
  }
  {\endlist}
  {\item}
% --------------------------------------
%            НАСТРОЙКИ СПИСКОВ
% --------------------------------------

\newlength{\listitemparagraphindent}
\setlength{\listitemparagraphindent}{1.25cm}
\newcommand{\listitemindent}{\hspace*{\listitemparagraphindent}}

\setlist{
	topsep=2pt,
	itemsep=2pt,
	parsep=0pt,
	partopsep=0pt
}

\setlist[itemize]{
    label=-,
    leftmargin=0pt,
    labelwidth=1.25cm,
    labelsep=0.5em,
    itemindent=1.8cm,
    align=right,
    listparindent=\listitemparagraphindent,
    parsep=0pt,
    itemsep=0pt,
    topsep=0pt
}

\setlist[enumerate]{
    leftmargin=0pt,
    labelwidth=1.6cm,
    labelsep=0.5em,
    itemindent=2cm,
    align=right,
    listparindent=\listitemparagraphindent,
    parsep=0pt,
    itemsep=0pt,
    topsep=0pt,
    label=\arabic*),
}

\setlist[enumerate,2]{
    leftmargin=0pt,
    labelwidth=1.6cm,
    labelsep=0.5em,
    itemindent=2.2cm,
    align=right,
    listparindent=\listitemparagraphindent,
    parsep=0pt,
    itemsep=0pt,
    topsep=0pt,
    label=\arabic*),
}

\hypersetup{
    pdftitle={Разработка системы защещенного хранения данных},
    pdfauthor={Куприянов Артём Максимович},
    pdfsubject={Дипломная работа},
    pdfkeywords={latex, диплом, безопасность},
    pdfcreator={LuaLaTeX},
    pdfproducer={TeX Live}
}

% --- Где искать картинки ---
\graphicspath{ {./figures/} }

% --- Основные настройки ---
\onehalfspacing % 1.5 отступ 
\setlength{\parindent}{1.25cm} % абзацный отступ
% Отступы после рисунков, таблиц, и.т.д
\setlength{\floatsep}{0pt}
\setlength{\intextsep}{0.8\baselineskip}
\setlength{\textfloatsep}{0.8\baselineskip}
\setlength{\belowcaptionskip}{0pt}
% Базовый отступ от заголовков
\setlength{\headsep}{0.5\baselineskip}

% Настройки приложений
\newcommand{\listingsappendix}{А}
\newcommand{\currentappendixletter}{\Asbuk{appendixcount}}
\newpagestyle{appendixcontstyle}{
    \sethead{}{Продолжение Приложения \currentappendixletter}{}
    \setfoot{}{\thepage}{}
}

% --- Подпись таблиц ---
\captionsetup[table]{
	justification=raggedright,
	singlelinecheck=false,
	font={stretch=1.5}
}

% --- Отступы для figure[p] ---

\makeatletter
\setlength{\@fptop}{0pt}
\setlength{\@fpsep}{\baselineskip}
\setlength{\@fpbot}{0pt plus 1fil}
\makeatother

% --- Запасной шрифт для кода ---
\directlua{
  luaotfload.add_fallback("monocyrfallback", {
    "PT Mono:mode=harf;script=cyrl;"
  })
}

% --- Команда, чтобы по максимуму вместить рисунок, если нужен именно на этой странице ---
\newcommand{\fitgraphics}[1]{\includegraphics[max width=\textwidth, max height=\pagegoal-\pagetotal-2cm, keepaspectratio]{#1}}
% --- Добавляет сноску вида "Рисунок 2, приложение Б, с. 54" ---
\newcommand{\listingfootnote}[1]{\footnote{Рисунок \ref{#1}, приложение \listingsappendix{}, с. \pageref{#1}}}
% --- Добавляет в содержание заголовок второго уровня с произвольным текстом ---
\newcommand{\addsubsubsection}[1]{\refstepcounter{subsubsection}\addcontentsline{toc}{subsubsection}{\protect\numberline{\thesubsubsection}#1}}

% --- Шрифты и переносы строк ---
\setmainfont{Times New Roman}
\setmonofont{Inconsolata}
\hypersetup{breaklinks=true}

\tolerance=2000
\emergencystretch=2em
\hyphenpenalty=1000
\exhyphenpenalty=1000
\hbadness=10000

\makeatletter
\lst@AddToHook{Init}{
  \tolerance=9999\relax
  \emergencystretch=0pt\relax
  \hyphenpenalty=0\relax
  \hbadness=10000\relax
}
\makeatother
% ------------------------------

\let\oldtexttt\texttt
\renewcommand{\texttt}[1]{%
  {\def\UrlFont{\ttfamily}\nolinkurl{#1}}%
}

% Отутспы от секций
% \titlespacing*{\section}{0pt}{0pt}{0.3cm}
% \titlespacing*{\subsection}{0pt}{0pt}{0pt}
% \titlespacing*{\subsubsection}{0pt}{0pt}{0pt}

% Заменяет разделитель между номером флоата (риснок, таблица) с двоеточия на дефис
\DeclareCaptionLabelSeparator{dash}{ -- }
\captionsetup{labelsep=dash}

\addto\captionsrussian{
    \renewcommand{\figurename}{Рисунок}
}

% Убирает автоматическую надпись "Содержание"
\addto\captionsrussian{\renewcommand{\refname}{}} 

%  --- Таблицы ---
\newcolumntype{L}[1]{>{\raggedright\arraybackslash}p{#1}}
\newcolumntype{Z}{>{\raggedright\arraybackslash}X}


\setlength{\LTpost}{0pt}

\makeatletter
\let\origtable\table
\let\endorigtable\endtable

\renewenvironment{table}
  {\setlength{\intextsep}{0cm}%
   \setlength{\textfloatsep}{0cm}%
   \origtable}
  {\endorigtable}
\makeatother
% -------------------

\begin{document}

	\includepdf[pages=-]{титульник.pdf}
	
	\includepdf[pages=-]{задание.pdf}

	\thispagestyle{empty}
	
	{\centering РЕФЕРАТ\par}
	
	\vspace*{\baselineskip}
	
	Дипломная работа содержит \pageref*{bib:end} с., \tablesBeforeBib{} \russianTables{\tablesBeforeBib}, \totalappendices{} \russianAppendices{\totalappendices}, \totalbibitems{} \russianSources{\totalbibitems}
	
	\vspace*{\baselineskip}
	
	КРИПТОГРАФИЯ, ШИФРОВАНИЕ, РАЗРАБОТКА, PYTHON, API
	
	\vspace*{\baselineskip}
	
	В работе разработана программа, централизующая процесс управления секретными данными, использующая передовые алгоритмы аутентифицированного шифрования.
	
	Целью данной дипломной работы является анализ существующих алгоритмов аутентифицированного шифрования, выбор подходящих в рамках данной работы, анализ существующих на рынке решений, решающих похожие задачи, составление технического задания на разработку, выбор подходящих технологий для разработки, анализ модели угроз, разработка ПО, и его тестирование.
	
	Актуальность данной работы заключается в необходимости управления секретными данными, такими как API-ключи, токены, пароли в окружении современных приложений. На данный момент большинство команд разработки используют подход переменных окружения для управления секретами, однако этот подход в определенных ситуациях представляет риски для утечки данных, а также при большом количестве секретных данных управление ими в файлах <<\texttt{.env}>> не всегда является удобным.
	
	Для решения этих проблем в рамках данной работы был разработан сервис, реализующий модель конвертного шифрования, централизующий управление секретными данными через интерфейс командной строки, и выдающий данные по веб-API.
		\newpage

	{\centering СОДЕРЖАНИЕ\par}
	
	\vspace*{-1.6cm}

	\tableofcontents
	\newpage
	
	\genericsection{Введение}
	
	По мере роста объёма цифровой информации и количества прикладных сервисов, важность безопасного хранения данных, а также возможности контроля доступа к этой информации, становится чрезвычайно высокой. Эффективные системы безопасности должны учитывать как криптографическую стойкость алгоритмов, так и надлежащее управление жизненным циклом ключей, проверку пользовательского ввода, контроль доступа и повторяемый операционный процесс. Особую озабоченность вызывает вопрос о том, как объединить требования этих разрозненных систем в единое программное решение, которое можно эффективно использовать на практике, а также масштабировать для будущих задач.
	
	В современных средах разработки команды одновременно должны поддерживать несколько сред: локальные, тестовые, эксплуатационные. Во всех них циркулируют чувствительные данные, такие как:

	\begin{itemize}
	\item ключи шифрования;
	\item токены API внешних сервисов;
	\item данные пользователей.
	\end{itemize}
	
	При отсутствии механизма управления секретнными данными возникают следующие проблемы:
	
	\begin{itemize}
	\item секреты хранятся в файлах <<\texttt{.env}>>, конфигурационных файлах, скриптах, что повышает риск утечки;
	\item ключи копируются вручную, из-за чего появляются проблемы с контролем версией ключей;
	\item процессы шифрования и дешифрования выполняются неконсистентно, что приводит к ошибкам интеграций;
	\item сложно отслеживать кто, когда, и при каких обстоятельствах получил доступ к данным.
	\end{itemize}
	
	При разработке это приводит к снижению скорости внедрения изменений и росту количества инцидентов при взаимодействии команд разработки и DevOps. Для сред эксплуатации это создает риски компрометации данных, простоя сервисов и нарушений требований регуляторов.
	
	В связи с этим возникает необходимость в решении, которое:
	
	\begin{itemize}
	\item стандартизирует операции с данными;
	\item позволяет унифицировать работу с данными в разных средах;
	\item обеспечивает удобный API для разработчиков.
	\end{itemize}

	Цель данной работы -- разработка программного решения для централизованного управления ключами, обеспечивающего безопасное хранение, повторное шифрование и предоставление зашифрованных данных.
		
	Для достижения указанной цели решаются следующие задачи:
	
	\begin{itemize}
		\item изучить современные алгоритмы шифрования и выбрать подходящие для данного проекта;
		\item проанализировать существующие решения на рынке;
		\item составить функциональные и нефункциональные требования к системе;
		\item спроектировать модель данных приложения;
		\item выбрать подходящий технологический стек для приложения;
		\item проанализировать модель угроз приложения;
		\item разработать программу;
		\item разработать модули тестирования программы.
	\end{itemize}
	
	Разработанный в рамках этой дипломной работы проект посвящен разработке программного приложения <<Vault>>, которое обеспечивает безопасное управление ключами и обработку данных. Приложение будет иметь интерфейс командной строки для администрирования, веб-API для предоставления данных после соответствующей авторизации, модуль криптографических операций и уровень хранения, использующий реляционную базу данных. Приложение использует современные алгоритмы AEAD (<<AES-128-GCM>>, <<AES-256-GCM>> и <<ChaCha20-Poly1305>>) для обеспечения конфиденциальности и целостности данных. 
	
	Одним из ключевых принципов проектирования является разделение ролей ключей. Например, главный ключ (KEK) никогда не хранится в базе данных, а предоставляется среде, в которой происходит операция, в виде зашифрованного (упакованного) значения. 
	
	При этом для шифрования данных используется другой ключ (ключ шифрования данных, DEK).
	
	Такая архитектура разделения ролей ключей ограничивает потенциальную уязвимость для компрометации как главного ключа, так и зашифрованных значений, хранящихся в базе данных, и соответствует основным принципам безопасных систем управления ключами.
	
	\newpage
	
	\section[Аналитический раздел]{АНАЛИТИЧЕСКИЙ РАЗДЕЛ}
	
	\subsection{Теоретические основы криптографии и модели безопасности}
	\subsubsection{Базовые свойства защиты (Триада CIA)}
	
	Фундамент теоретических представлений об информационной безопасности образует так называемая <<триада CIA>> -- конфиденциальность, целостность и доступность (Confidentiality, Integrity, Availability). Данные принципы имеют междисциплинарный характер и составляют основу как законодательной, так и технической областей \cite{milnichuk}.
	
	\begin{itemize}
	\item \textbf{confidentiality (конфиденциальность)} означает, что информация не может быть прочитана или использована лицами, не имеющими соответствующих полномочий. Обеспечение
	конфиденциальности предполагает решение двуединой задачи: предотвращение несанкционированного доступа к данным и недопущение их перехвата при передаче.
	Практическая реализация этого принципа осуществляется посредством шифрования, разграничения доступа, физической защиты носителей \cite{milnichuk};
	\item \textbf{integrity (целостность)} заключается в ее свойстве существовать в неискаженном виде и не подвергаться случайным или преднамеренным изменениям. Нарушение целостности может быть следствием как злонамеренных действий (внедрение вредоносного
	ПО, модификация данных), так и технических сбоев. Защита целостности требует применения методов контроля (контрольные суммы, хэш-функции) и предотвращения (резервирование, управление доступом) \cite{milnichuk};
	\item \textbf{availability (доступность)} предполагает, что авторизованные пользователи получают доступ к информации и связанным с ней активам всегда, когда это необходимо. Угрозы доступности (например, DDoS-атаки) способны парализовать деятельность организаций, критически зависимых от информационных систем. Обеспечение доступности достигается резервированием компонентов, созданием отказоустойчивых конфигураций, планированием восстановления после сбоев \cite{milnichuk}.
	\end{itemize}
	
	
	Теоретический анализ информационной безопасности неразрывно связан с
	классификацией угроз. Под угрозой информационной безопасности понимается совокупность условий и факторов, создающих потенциальную или реальную опасность нарушения безопасности информации \cite{milnichuk}.
		
	\begin{figure}[H]
%	\fitgraphics{cia.pdf}
	\centering\includegraphics[width=15cm]{cia.pdf}
	\caption{Взаимодействие компонентов триады CIA}
	\end{figure}
	
	\subsubsection{Используемые криптографические термины}

%	\addsubsubsection{Алгоритмы аутентифицированного шифрования}
	\textbf{Алгоритмы аутентифицированного шифрования} (Authenticated Encryption -- АЕ) \label{def:AEAD}
	представляют собой класс алгоритмов шифрования, которые также позволяют аутентифицировать зашифрованные данные, т. е. обеспечивают возможность проверки их
	целостности.
	Алгоритмы аутентифицированного шифрования с присоединёнными данными
	(Authenticated Encryption with Associated Data -- AEAD) можно считать частным случаем АЕ-алгоритмов, поскольку они также обеспечивают как конфиденциальность,
	так и целостность зашифрованных данных. Помимо этого, AEAD-алгоритмы позволяют обеспечить целостность и аутентичность некоторых дополнительных данных,
	для которых не требуется конфиденциальность. Дополнительные данные при этом не
	шифруются.
	Необходимость в появлении такого класса алгоритмов обусловлена тем, что во многих приложениях требуется, помимо передачи некоторых данных в зашифрованном
	виде, передавать также определённую открытую информацию с контролем её целостности/аутентичности. Данная возможность важна, когда взаимодействие с устройствами с небольшими ресурсами ограничено незначительным промежутком времени.
	
	Разработанный проект использует алгоритмы симметричного шифрования AEAD <<AES-128-GCM>>, <<AES-256-GCM>> и <<ChaCha20-Poly1305>> для проверки целостности данных.	
	
%	\addsubsubsection{Ключ шифрования ключа (Key Encryption Key, KEK)}
	\textbf{Ключ шифрования ключа (Key Encryption Key, KEK)} -- ключ, используемый для шифрования другого ключа. (обычно ключей шифрования трафика -- Traffic Encryption Key -- TEK) для передачи или хранения \cite{CNSSI_4009_2015}.
	
%	\addsubsubsection{Ключ шифрования данных (Data Encryption Key, DEK)}
	\textbf{Ключ шифрования данных (Data Encryption Key, DEK)} -- Ключ, используемый для шифрования и расшифрования данных, не являющихся другими ключами \cite{NIST_SP_800-57_Part_1_Rev._5}.
%	AEAD-алгоритмы позволяют за максимально возможно короткий промежуток времени (ограниченный однократным выполнением алгоритма) сформировать ответ, содержащий следующее:
%	 
%	\begin{itemize}
%	\item зашифрованную полезную информацию
%	данные
%	\item присоединённые данные
%	\item код аутентификации полученных данных — тег, позволяющий проверить целостность и присоединённых данных, и полезной информации. \cite{panasenko}
%	\end{itemize}

%	\addsubsubsection{Однократно используемое число (nonce)}
	\textbf{Однократно используемое число (nonce)} -- Значение, используемое в криптографических протоколах, которое никогда не повторяется с одним и тем же ключом \cite{GOST_R_34.14.2025}.
	
%	\addsubsubsection{Имитовствка (tag)}
	\textbf{Имитовставка (tag)} -- Битовая строка, добавляемая к сообщению с целью обеспечения возможности обнаружения подмены и/или имитации сообщения и являющаяся результатом применения к нему криптографической хеш-функции, зависящей от ключа \cite{GOST_R_34.14.2025}.
	
	\textbf{Оборачивание ключа (key wrapping)} -- это криптографическая операция шифрования ключевого материала с использованием ключа шифрования ключей (KEK), направленная на обеспечение конфиденциальности и целостности защищаемого ключа при его хранении или передаче.
	
	\textbf{Извлечение ключа (key unwrapping)} -- это операция, обратная оборачиванию, при которой исходный ключ восстанавливается из защищенного (обернутого) представления с использованием ключа шифрования ключей (KEK).
		
	\subsubsection{Общие термины, используемые в работе}
	
	\textbf{Dev (development)} -- среда, в которой выполняется разработка, тестирование и отладка программного обеспечения.
	
	\textbf{Prod (production)} -- промышленная среда эксплуатации, в которой система используется для обработки реальных данных и предъявляет повышенные требования к надёжности и безопасности.
	
	Среда \textbf{staging (предпродуктивная среда)} представляет собой изолированную среду, максимально приближенную к <<production>>, предназначенную для проверки корректности работы системы перед развертыванием в промышленной эксплуатации. В данной среде проводится интеграционное тестирование, проверка конфигурации и сценариев развертывания.
		
	\subsubsection{Алгоритм PBKDF2}	
	
	\label{def:PBKDF2}
	
	В разработанной программе для вывода ключей используется криптографическая функция PBKDF2.
	
	Алгоритм вывода ключей PBKDF2 (Password-Based Key Derivation Function 2) необходим для вывода стойкого криптографически ключа из строки пользователя. Алгоритм определен в стандарте RFC 8018 \cite{RFC8018} и основан на применении PRF (псевдослучайной функции). 
	
	Суть алгоритма заключается в увеличении сложности перебора (Brute Force) паролей за счет применения соли и большого числа итераций. Схема работы алгоритма PBKDF2 представлена на рисунке.
	
	\begin{figure}[h]
%		\centering\fitgraphics{PBKDF2.pdf}
		\centering\includegraphics[width=10cm]{PBKDF2.pdf}
		\caption{Схема работы алгоритма PBKDF2}
		\label{fig:PBKDF2}
	\end{figure}
	
	Алгоритм принимает на вход следующие параметры:
	
	\begin{itemize}
	\item пароль $P$;
	\item соль $S$;
	\item число итераций $c$;
	\item длину ключа $dkLen$.
	\end{itemize}
	
	В основе PBKDF2 лежит псевдослучайная функция, обычно выбирается HMAC.
	
	\begin{enumerate}
	
	\item Начальное вычисление
	
	Сначала выбирается значение:
	
	\begin{equation}
		U_{1} = PRF(P,S||i)
	\end{equation}
	
	где 
	
	\begin{itemize}
	\item $i$ -- индекс блока выходного ключа;
	\item $||$ -- операция сложения строк (конкатенация).
	\end{itemize}
	
	\item Итеративное преобразование
	
	Далее выполняется последовательное вычисление значений:
	
	\begin{equation}
		U_{j} = PRF(P,U_{j-1}), j = 2,3,...,c
	\end{equation}
	
	Таким образом, результат предыдущей итерации используется как вход для следующей, что увеличивает вычислительную сложность алгоритма.
	
	\item Агрегация результатов
	
	Итоговое значение функции $F$ вычисляется как побитовое сложение по модулю 2 (XOR) всех промежуточных значений:
	
	\begin{equation}
		F(P,S,c,i) = U_1 \oplus U_2 \oplus \cdots \oplus U_c
	\end{equation}
	
	\item Формирование итогового ключа
	
	Выходной ключ формируется путём конкатенации результатов функции $F$ для различных значений индекса $i$:
	
	\begin{equation}
		DK = PBKDF2(P,S,c,dkLen)
	\end{equation}
	
	\end{enumerate}
	
	\subsubsection{Алгоритм AES-GCM}
	
	\label{def:AESGCM}
	
	В разработанной системе для шифрования данных используются алгоритмы семейства AES-GCM.
	
	AES-GCM (Advanced Encryption Standard в режиме Galois/Counter Mode) представляет собой симметричный алгоритм аутентифицированного шифрования (AEAD), обеспечивающий одновременно конфиденциальность и контроль целостности данных. Алгоритм стандартизирован в NIST SP 800-38D \cite{NISTGCM}.
	
	Суть алгоритма заключается в комбинировании режима потокового шифрования (CTR) и вычисления аутентификационного тега с использованием функции GHASH. Схема работы алгоритма AES-GCM представлена на рисунке \ref{fig:AESGCM}.
	
	\begin{figure}[h]
		\centering\includegraphics[width=15cm]{AESGCM.pdf}
		\caption{Схема работы алгоритма AES-GCM}
		\label{fig:AESGCM}
	\end{figure}
	
	Алгоритм принимает на вход следующие параметры:
	
	\begin{itemize}
	\item секретный ключ $K$;
	\item nonce (инициализационный вектор) $N$;
	\item открытый текст $P$;
	\item дополнительные аутентифицируемые данные $A$.
	\end{itemize}
	
	В результате работы алгоритма формируются:
	\begin{itemize}
	\item зашифрованный текст $C$;
	\item тег аутентификации $T$.
	\end{itemize}
	
	\begin{enumerate}

	\item Формирование вспомогательного ключа
	
	Сначала вычисляется вспомогательный ключ для функции GHASH:
	
	\begin{equation}
		H = AES_K(0^{128})
	\end{equation}
	
	где $H$ используется как ключ хеширования в поле $GF(2^{128})$.
	
	\item Шифрование данных
	
	Шифрование выполняется в режиме счетчика (CTR):
	
	\begin{equation}
		C = AES\text{-}CTR_K(P, N)
	\end{equation}
	
	В данном режиме генерируется поток ключа, который складывается по модулю 2 с открытым текстом.
	
	\item Вычисление тега аутентификации
	
	Для обеспечения целостности данных вычисляется значение GHASH:
	
	\begin{equation}
		S = GHASH_H(A, C)
	\end{equation}
	
	где учитываются как зашифрованные данные, так и дополнительные аутентифицируемые данные.
	
	\item Формирование итогового тега
	
	Итоговый тег аутентификации определяется как:
	
	\begin{equation}
		T = S \oplus AES_K(N)
	\end{equation}
	
	\item Результат работы алгоритма
	
	Итогом работы AES-GCM является пара значений:
	
	\begin{equation}
		(C, T) = AES\text{-}GCM_K(P, N, A)
	\end{equation}
	
	\end{enumerate}
	
	\subsubsection{Алгоритм ChaCha20-Poly1305}
	
	\label{def:ChaChaPoly}
	
	В разработанной системе наряду с AES-GCM используется алгоритм ChaCha20-Poly1305.
	
	ChaCha20-Poly1305 представляет собой алгоритм аутентифицированного шифрования (AEAD), сочетающий потоковый шифр ChaCha20 и механизм аутентификации сообщений Poly1305. Алгоритм стандартизирован в RFC 8439 \cite{RFC8439}.
	
	Основная идея алгоритма заключается в использовании ChaCha20 для шифрования данных и Poly1305 для вычисления тега аутентификации. Схема работы алгоритма представлена на рисунке \ref{fig:ChaChaPoly}.
	
	\begin{figure}[h]
		\centering\includegraphics[width=15cm]{ChaChaPoly.pdf}
		\caption{Схема работы алгоритма ChaCha20-Poly1305}
		\label{fig:ChaChaPoly}
	\end{figure}
	
	Алгоритм принимает на вход следующие параметры:
	
	\begin{itemize}
	\item секретный ключ $K$;
	\item nonce $N$;
	\item открытый текст $P$;
	\item дополнительные аутентифицируемые данные $A$.
	\end{itemize}
	
	В результате работы формируются:
	\begin{itemize}
	\item зашифрованный текст $C$;
	\item тег аутентификации $T$.	
	\end{itemize}
	
	\begin{enumerate}
	
	\item Генерация ключа Poly1305
	
	На первом этапе с использованием ChaCha20 генерируется одноразовый ключ для алгоритма Poly1305:
	
	\begin{equation}
		(r, s) = ChaCha20_K(N, 0)
	\end{equation}
	
	где первый блок выходного потока используется как ключ аутентификации.
	
	\item Шифрование данных
	
	Шифрование выполняется с использованием потокового шифра ChaCha20:
	
	\begin{equation}
		C = ChaCha20_K(P, N)
	\end{equation}
	
	Открытый текст складывается по модулю 2 с псевдослучайным потоком ключа.
	
	\item Формирование данных для аутентификации
	
	Для вычисления тега формируется последовательность:
	
	\begin{equation}
		M = A \parallel C
	\end{equation}
	
	где $||$ обозначает операцию конкатенации.
	
	\item Вычисление тега аутентификации
	
	Тег вычисляется с использованием алгоритма Poly1305:
	
	\begin{equation}
		T = Poly1305_{r}(M)
	\end{equation}
	
	\item Результат работы алгоритма
	
	Итогом работы является пара значений:
	
	\begin{equation}
		(C, T) = ChaCha20\text{-}Poly1305_K(P, N, A)
	\end{equation}
	
	\end{enumerate}
	
	\subsection{Анализ предметной области}
			
	\subsubsection{Особенности задачи для сред разработки и production}
	
	Основная проблема при работе с чувствительными данными заключается в том, что требования безопасности часто внедряются во вред удобству.
	
	Для среды разработки важны:
	
	\begin{itemize}
	\item простота локального запуска и тестирования;
	\item понятный CLI интерфейс для типовых операций;
	\item быстрая диагностика ошибок в формате, удобном для CI/CD.
	\end{itemize}
	
	Для production-среды важны:
	
	\begin{itemize}
	\item недоступность мастер-ключа в базе данных и репозитории;
	\item шифрование с контролем целостности (AEAD);
	\item защищенный сетевой доступ (HTTPS + авторизация);
	\item контролируемая схема хранения с миграционной совместимостью.
	\end{itemize}
	
	Данный проект решает этот конфликт через разделение ответственности компонентов: CLI используется для администрирования и загрузки данных, API -- для строго ограниченной выдачи данных. При этом криптографическая логика и валидация вынесены в отдельные модули и не дублируются между сценариями.
	
	\subsubsection{Анализ существующих решений (конкурентов)}
	
	На данный момент на рынке уже существуют зрелые системы управления чувствительными данными. Их использование оправдано в крупных инфраструктурах, однако для учебных, исследовательских и части прикладных задач они могут быть избыточны по сложности и стоимости внедрения.
	

	\begin{enumerate}                                        

	\item HashiCorp Vault
	
	HashiCorp Vault -- это кроссплатформенное программное обеспечение с открытым исходным кодом, предназначенное для централизованного хранения, защиты и управления доступом к чувствительным данным (секретам), таким как пароли, API-ключи, сертификаты и токены.
	
	Сильные стороны:
	
	\begin{itemize}
	\item широкий набор механизмов аутентификации и авторизации;
	\item динамические секреты и развитые политики доступа;
	\item интеграции с Kubernetes, облачными провайдерами и PKI.
	\end{itemize}
	
	Ограничения для данной задачи:
	
	\begin{itemize}
	\item существенная операционная сложность (инициализация, unseal, сопровождение кластера);
	\item более высокий порог входа для учебного проекта и небольших команд.
	\end{itemize}
		
	\item Облачные сервисы (AWS KMS/Secrets Manager, Azure Key Vault, Google Cloud KMS/Secret Manager)
	
	В современной практике разработки выделяют две основные группы облачных инструментов: KMS (Key Management Service) для управления ключами шифрования и Secret Manager для хранения паролей и токенов.
	
	\begin{enumerate}[label=\arabic{enumi}.\arabic*), ref=\arabic{enumi}.\arabic*]
	\item AWS (Amazon Web Services)
	\begin{itemize}
	\item AWS KMS -- Централизованное управление ключами (CMK). Использует аппаратные модули безопасности (HSM) уровня FIPS 140-2. Основная задача -- шифрование дисков (EBS), баз данных (RDS) и объектов (S3);
	\item AWS Secrets Manager -- Сервис для хранения секретов. Главная функциональность -- нативная ротация паролей для баз данных AWS (RDS, Redshift) без участия разработчика.
	\end{itemize}
	\item Microsoft Azure
	\begin{itemize}
		\item Azure Key Vault -- универсальный сервис <<три в одном>>. Позволяет хранить:
		\begin{itemize}[label=--]
		\item Secrets (пароли, строки подключения);
		\item Keys (ключи для шифрования дисков и данных);
		\item Certificates (управление SSL/TLS сертификатами).
		\end{itemize}
		Интегрируется с Azure Managed Identities, что позволяет приложениям получать доступ к секретам без использования паролей.
	\end{itemize}
	\item Google Cloud Platform (GCP)
	\begin{itemize}
	\item Cloud KMS -- Позволяет управлять симметричными и асимметричными ключами. Поддерживает внешние менеджеры ключей (EKM), если данные должны шифроваться ключами вне облака Google;
	\item Secret Manager -- Глобальный сервис для хранения конфиденциальных данных с удобным версионированием секретов и разделением доступа через Cloud IAM.
	\end{itemize}
	\end{enumerate}
	
	Сильные стороны облачных сервисов:
	
	\begin{itemize}
	\item высокая надежность;
	\item интеграция с конкретными облачными сервисами;
	\item встроенное логирование и IAM-модели.
	\end{itemize}
	
	Ограничения для данной задачи:
	
	\begin{itemize}
	\item привязка к конкретному облачному провайдеру;
	\item зависимость от внешней инфраструктуры.
	\end{itemize}
	
	\item Корпоративные PAM/Secrets-решения (например, CyberArk, Delinea)
	
	PAM (Privileged Access Management) -- это комплекс стратегий и технологий для контроля, мониторинга и аудита привилегированного доступа к критическим ресурсам организации. В отличие от легковесных инструментов для DevOps, PAM-системы исторически сфокусированы на управлении доступом администраторов и операторов (Human-to-Machine).

	Сильные стороны:
	
	\begin{itemize}
	\item интеграция с корпоративной политикой информационной безопасности;
	\item глубокие процессы аудита и контроля.
	\end{itemize}
	
	Ограничения для данной задачи:
	
	\begin{itemize}
	\item высокая стоимость;
	\item избыточность в малых и средних компаниях.
	\end{itemize}
	
	\item Open-source self-hosted решения (например, OpenBao, Infisical self-hosted)
	
	Сильные стороны:
	
	\begin{itemize}
	\item контроль над инфраструктурой и отсутствие жесткой привязки к конкретному облаку;
	\item более низкий порог по стоимости владения по сравнению с enterprise PAM;
	\item гибкость кастомизации под внутренние процессы команды.
	\end{itemize}
	
	Ограничения для данной задачи:
	
	\begin{itemize}
	\item ответственность за эксплуатацию (обновления, бэкапы, мониторинг) полностью на команде;
	\item часто менее зрелая экосистема интеграций по сравнению с лидерами рынка.
	\end{itemize}
	
	\item SaaS-сервисы для разработчиков (например, Doppler, 1Password Secrets Automation)
	
	Это полностью облачные платформы, которые берут на себя всю сложность управления инфраструктурой (шифрование, репликацию, бэкапы), предоставляя разработчикам только интерфейс и API. Основной акцент сделан на Developer Experience (DX) -- максимальную простоту интеграции секретов в код.
	
	Сильные стороны:
	
	\begin{itemize}
	\item быстрый старт и удобная интеграция в CI/CD;
	\item низкая операционная нагрузка на команду;
	\item хороший UX для небольших и средних продуктовых команд.
	\end{itemize}
	
	Ограничения для данной задачи:
	
	\begin{itemize}
	\item зависимость от внешнего провайдера и его модели безопасности;
	\item ограниченная гибкость в нетиповых сценариях и в локальных окружениях;
	\item возможный рост затрат при масштабировании.
	\end{itemize}
	
	\item GitOps-инструменты шифрования конфигураций (например, Mozilla SOPS, git-crypt, Sealed Secrets)
	
	Сильные стороны:
	
	\begin{itemize}
	\item удобная работа с секретами в Git-репозиториях и workflow инфраструктурного кода;
	\item простая интеграция в существующие процессы поставки конфигураций;
	\item низкий порог внедрения для задач <<шифровать и хранить в репозитории>>.
	\end{itemize}
	
	Ограничения для данной задачи:
	
	\begin{itemize}
	\item это не полноценные платформы управления жизненным циклом секретов;
	\item ограниченные возможности централизованной выдачи секретов по API и детального policy-management;
	\item при росте масштаба требуется дополнительная инфраструктура и процессы контроля.
	\end{itemize}
	
	\end{enumerate}
		
	\subsubsection{Выводы по анализу и обоснование выбора собственного решения}
	
	Результаты проведенного анализа схематично показаны на рисунке \ref{fig:solutions_quadrant}.
	
	Также сравнительный анализ решений представлен в таблице \ref{tab:solutions_compare_full}.

%	\Needspace{0.55\textheight}
	{\fontsize{12pt}{14pt}\selectfont
	\begin{tabularx}{\textwidth}{|p{3.8cm}|Z|Z|}
\caption{Сравнительный анализ существующих решений для управления секретами}
\label{tab:solutions_compare_full} \\

\hline
\multicolumn{1}{|c|}{Тип решения} &
\multicolumn{1}{c|}{Сильные стороны} &
\multicolumn{1}{c|}{Ограничения для данной задачи} \\ \hline
\multicolumn{1}{|c|}{1} &
\multicolumn{1}{c|}{2} &
\multicolumn{1}{c|}{3} \\ \hline
\endfirsthead

%\multicolumn{3}{r}{} \\[-0.3em]
\multicolumn{3}{r}{{\fontsize{14pt}{14pt}\selectfont Продолжение таблицы \thetable}} \\[0.5em]
\hline
%\multicolumn{1}{|c|}{Тип решения} &
%\multicolumn{1}{c|}{Сильные стороны} &
%\multicolumn{1}{c|}{Ограничения для данной задачи} \\ \hline
\multicolumn{1}{|c|}{1} &
\multicolumn{1}{c|}{2} &
\multicolumn{1}{c|}{3} \\ \hline
\endhead

SaaS-сервисы (Doppler, 1Password Secrets Automation) &
Быстрый старт; удобная интеграция в CI/CD; низкая операционная нагрузка; хороший UX &
Зависимость от провайдера; ограниченная гибкость; возможный рост затрат \\
\hline

HashiCorp Vault &
Широкий набор механизмов аутентификации и авторизации; динамические секреты; развитые политики доступа; интеграции с Kubernetes, облачными провайдерами и PKI &
Существенная операционная сложность; более высокий порог входа для учебного проекта и небольших команд \\
\hline

Облачные сервисы (AWS KMS / Secrets Manager, Azure Key Vault, Google Cloud KMS / Secret Manager) &
Высокая надёжность; интеграция с конкретными облачными сервисами; встроенное логирование и IAM-модели &
Привязка к конкретному облачному провайдеру; зависимость от внешней инфраструктуры \\
\hline

Корпоративные PAM/Secrets-решения (CyberArk, Delinea) &
Интеграция с корпоративной политикой информационной безопасности; глубокие процессы аудита и контроля &
Высокая стоимость; избыточность для малых и средних компаний и учебных задач \\
\hline

Open-source self-hosted решения (OpenBao, Infisical self-hosted) &
Контроль над инфраструктурой; отсутствие привязки к облаку; более низкая стоимость владения; гибкость кастомизации &
Ответственность за эксплуатацию полностью лежит на команде; менее зрелая экосистема интеграций \\
\hline

GitOps-инструменты (SOPS, git-crypt, Sealed Secrets) &
Удобная работа с секретами в Git; простая интеграция; низкий порог внедрения &
Не являются полноценными системами управления секретами; ограниченный lifecycle management \\
\hline

\end{tabularx}
	}

	\begin{figure}[!hbt]
%		\centering\fitgraphics{quadrant\_choice.pdf}
		\centering\includegraphics[width=10cm]{quadrant\_choice.pdf}
		\caption{Сравнение решений для управления секретами}
		\label{fig:solutions_quadrant}
	\end{figure}	
	
	Выполненный анализ показал, что существующие на данный момент платформы хорошо решают задачу управления секретами в промышленном масштабе, но не оптимальны для сценария, в котором необходимы:
		
	\begin{itemize}
	\item прозрачная архитектура;
	\item контролируемая сложность внедрения;
	\item отсутствие зависимости от внешних облачных сервисов;
	\item единый поток работы для <<dev>> и <<prod>> с воспроизводимыми командами.
	\end{itemize}	
	
	Именно поэтому в данной работе выбран подход создания специализированного сервиса. Он решает ключевые практические потребности: безопасное хранение DEK в обернутом виде, использование KEK вне базы данных, перешифрование данных между алгоритмами, авторизованный доступ по API и тестируемость всех критических операций. Такой подход обеспечивает баланс между реализуемостью и безопасностью. 	
		
	\FloatBarrier	
	
	\clearpage
	
	\section[Практический раздел]{ПРАКТИЧЕСКИЙ РАЗДЕЛ}
	
	\subsection{Постановка задачи и требования к системе}
	
	Назначение разрабатываемой системы заключается в реализации безопасного процесса работы с криптографическими ключами и данными в средах разработки и <<production>>. Постановка задачи состоит из следующих частей:
	
	\begin{itemize}
		\item функциональные требования;
		\item нефункциональные требования;
		\item ограничения реализации;
		\item критерии приемки.
	\end{itemize}
	
	Функциональные требования определяют минимально необходимый набор возможностей системы. Они описывают, какие операции пользователь и администратор должны выполнять в рамках единого эксплуатационного контура.
	
	\begin{itemize}
	\item управление ключами через CLI. Система должна поддерживать генерацию и импорт ключевого материала командами \texttt{generate-kek, generate-dek, add-kek} и \texttt{add-dek}. Это необходимо для того, чтобы администратор мог как создавать новые ключи внутри системы, так и подключать заранее подготовленные значения. Одновременно интерфейс обязан валидировать формат данных, допустимые алгоритмы и логические имена, чтобы в хранилище не попадали некорректные или двусмысленные записи;
	\item безопасная загрузка и перешифрование данных. Команда \texttt{load} должна принимать входной контейнер, выполнять его расшифрование с использованием исходного DEK и затем повторное шифрование в целевой алгоритм с использованием выходного DEK. Такое требование отражает прикладной сценарий миграции данных между ключами и алгоритмами, где важно не просто хранить шифртекст, а контролируемо преобразовывать его внутри доверенного процесса;
	\item хранение данных и метаданных в БД. После выполнения криптографических операций система должна сохранять не только полезную нагрузку, но и достаточный набор служебных данных: имена ключей, алгоритмы, заголовок контейнера, временные метки и авторизационный секрет в защищенном виде. Это делает каждую запись самодостаточной для последующего чтения через API, но не приводит к хранению открытых секретов;
	\item выдача данных по API. Эндпоинт \texttt{GET /data/{name}} должен возвращать расшифрованные данные только после успешной проверки Bearer-авторизации, корректного чтения ключевой записи и успешного извлечения целевого DEK. Тем самым система выступает не только как инструмент подготовки контейнеров, но и как прикладной сервис контролируемого доступа к данным;
	\item сохранение нового открытого текста без вспомогательных скриптов. Поскольку в реальном эксплуатационном контуре возникают не только задачи перешифрования, но и задачи первичного ввода данных, система должна поддерживать сценарий добавления новых открытых данных. Для этого реализуется команда \texttt{store-plaintext}, которая читает открытый текст из файла, автоматически шифрует его выбранным алгоритмом и сохраняет результат в \texttt{data\_records}. Это требование закрывает практический пробел между криптографическим API и пользовательским сценарием административной загрузки данных.
	\end{itemize}
	
	Нефункциональные требования описывают следующие характеристики:
	
	\begin{itemize}
		\item безопасность. Это базовое нефункциональное требование для всего проекта. Система должна использовать современные AEAD-алгоритмы (<<AES-128-GCM>>, <<AES-256-GCM>>, <<ChaCha20-Poly1305>>), чтобы одновременно обеспечивать конфиденциальность и контроль целостности данных. Дополнительно требуется строгое разделение ролей KEK и DEK: master KEK не должен сохраняться в БД, а DEK должны храниться только в обернутом виде. Такое требование снижает последствия компрометации хранилища и соответствует принятой модели конвертного шифрования;
		\item легкость тестирования. Поведение системы должно быть формализуемым и воспроизводимо проверяемым через unit-, интеграционные и сценарные тесты. Для защищенного сервиса этого недостаточно описать словами: ключевые инварианты, такие как запрет хранения открытых DEK, отказ при неверном Bearer-токене, проверка заголовка контейнера и обязательность master KEK, должны подтверждаться автоматическими проверками;
		\item воспроизводимость. Сервис должен одинаково работать в dev- и production-подобных средах при различии только параметров конфигурации, а не программной логики. Это означает единый набор CLI-команд, единые форматы контейнеров, единое поведение API и предсказуемую работу через переменные окружения;
		\item переносимость. Решение не должно быть жестко привязано к конкретному облачному провайдеру, внешнему KMS или специализированной платформе. Такой подход важен для учебно-прикладного проекта, поскольку позволяет запускать проект локально, в лаборатории или на выделенном сервере без внешней инфраструктурной зависимости.
	\end{itemize}
	
	Ограничения и допущения описывают рамки текущей версии разработки:
	
	\begin{itemize}
	\item Master KEK хранится только в env-переменной;
	\item открытый KEK не хранится в базе данных и не фиксируется в репозитории;
	\item доступ к данным авторизуется через Bearer-токен, который хранится в виде PBKDF2-хеша\footnote{См. раздел с описанием работы алгоритма PBKDF2 на с. \pageref{def:PBKDF2}};
	\item в качестве хранилища используется реляционная база данных.
	\end{itemize}
	
	Критерии приемки обозначают условия, при которых задача считается выполненной:
	
	\begin{itemize}
	\item криптографические преобразования выполняются корректно;
	\item CLI-сценарии завершаются успешно;
	\item API-сценарии завершаются успешно;
	\item автотесты проекта выполняются успешно.
	\end{itemize}
	
	\subsection{Проектирование системы и модели данных}
	
	\subsubsection{Модули приложения}
		
	Проектирование системы выполнено в модульно-слоистой архитектуре, где каждый слой решает ограниченный набор задач и имеет явные точки взаимодействия. Распределение модулей по семантическим слоям представлено в таблице~\ref{tab:semantic}.
	
	{\fontsize{12pt}{14pt}\selectfont
    \begin{tabularx}{\textwidth}{|l|>{\raggedright\arraybackslash}X|}
	\caption{Распределение модулей по семантическим слоям}
	\label{tab:semantic} \\
	\hline
	\multicolumn{1}{|c|}{Семантическое значение} &
	\multicolumn{1}{c|}{Модуль} \\
%	\hline
%	\multicolumn{1}{|c|}{1} &
%	\multicolumn{1}{c|}{2} \\			
	\hline
	\endfirsthead
	
	
	\multicolumn{2}{r}{{\fontsize{14pt}{14pt}\selectfont Продолжение таблицы \thetable}} \\[0.5em]
	\hline
	\multicolumn{1}{|c|}{1} &
	\multicolumn{1}{c|}{2}
	\endhead			
	
	Точка входа & \texttt{main.py, vault/cli.py, vault/api.py} \\
	\hline
	Сервисный слой & \texttt{vault/crypto.py, vault/key\_manager.py, vault/auth.py} \\
	\hline
	Слой доступа к данным & \texttt{vault/db.py} (репозиторий и ORM-модели) \\
	\hline
	Слой валидации контрактов & \texttt{vault/schemas.py, vault/db\_schemas.py} \\
	\hline
	Конфигурационный слой & \texttt{vault/settings.py} и переменные окружения \\
	\hline
\end{tabularx}
    }

	Такое разделение снижает связанность кода: CLI и API используют общие криптографические и валидационные модули, а работа с базой данных изолирована в репозитории \texttt{VaultRepository}.
	
	Схема модулей приложения представлена на рисунке \ref{fig:modules}.

	\begin{figure}[!hbt]
%		\centering\includegraphics[width=\textwidth]{modules_alt.pdf}
%		\centering\fitgraphics{modules_alt.pdf}
		\centering\includegraphics[height=0.65\textheight]{modules_alt_subgraphs.pdf}
%		\centering\fitgraphics{modules_alt_subgraphs.pdf}
		\caption{Схема модулей приложения}
		\label{fig:modules}
	\end{figure}

	\subsubsection{Модель данных}
	
	Модель данных построена на двух ключевых сущностях: \texttt{keys}\footnote{См. определение таблицы \texttt{key\_records}: метка \calloutfullreffootnotetextonly{lst:table:keyrecord}} и \texttt{data\_records}\footnote{См. определение таблицы \texttt{data\_records}: метка \calloutfullreffootnotetextonly{lst:table:datarecord}}. Диаграмма сущностей представлена на рисунке \ref{fig:keys_and_data_records}.
	
	\begin{figure}[h]
	\centering\includegraphics[width=10cm]{keys\_data\_records.pdf}
	\caption{Таблицы \texttt{keys} и \texttt{data\_records}}
	\label{fig:keys_and_data_records}
	\end{figure}	

	Таблица \texttt{keys} хранит ключевой материал и метаданные ключей:
	
	\begin{enumerate}
	\item Ограничение уникальности \texttt{uq\_keys\_name\_type} не допускает дубли пары (\texttt{name, key\_type}).
	\item Для \texttt{key\_type="kek"} поле \texttt{key\_hex} хранит ключ в шестнадцатеричном представлении.
	\item Для \texttt{key\_type="dek"} поле \texttt{key\_hex} хранит обернутое представление DEK.
	\end{enumerate}
	
	Таблица \texttt{data\_records} хранит зашифрованные полезные данные и параметры их обработки:
		
	\begin{enumerate}
	\item Ограничение уникальности \texttt{uq\_data\_records\_name} обеспечивает upsert-семантику по имени записи.
	\item Поле \texttt{auth\_hash} хранит PBKDF2-хеш, а не исходный заголовок авторизации.
	\item Поля \texttt{alg\_in} и \texttt{alg\_out} фиксируют входной и выходной алгоритмы в процессе перешифрования.
	\item Поля \texttt{dek\_in} и \texttt{dek\_out} содержат ссылки на логические имена DEK, используемые в \texttt{load} и \texttt{GET /data/{name}}.
	\item Поля \texttt{header\_json} и \texttt{payload\_hex} хранят контейнер шифртекста: заголовок AEAD (alg/nonce/tag) и шифртекст в шестнадцатеричном представлении.
	\end{enumerate}
	
	На уровне контрактов данных применяются строгие ограничения:
	
	\begin{enumerate}
	\item Имена сущностей ограничены по длине (1..128), алгоритмы (1..64), заголовок авторизации (1..512)
	\item \texttt{payload\_hex} валидируется как корректная шестнадцатеричная строка
	\item \texttt{header\_json} проверяется как JSON-объект с ожидаемой структурой
	\item Обернутый DEK валидируется как JSON с полями \texttt{nonce, wrapped, tag}
	\end{enumerate}
	
	\subsubsection{Потоки обработки}
	
	Ключевые потоки обработки строятся вокруг двух сценариев.
	
	
	\begin{itemize}
	
	\item сценарий \texttt{load}:
	
	\begin{enumerate}
	\item Валидация входных параметров команды и формата контейнера.
	\item Получение \texttt{dek\_in} и \texttt{dek\_out} из таблицы \texttt{keys}.
	\item Извлечение DEK через master KEK из \texttt{VAULT\_MASTER\_KEK\_HEX}.
	\item Дешифрование входных данных (\texttt{alg\_in}) и повторное шифрование (\texttt{alg\_out}).
	\item Upsert записи в \texttt{data\_records} с обновлением \texttt{auth\_hash}, \texttt{header\_json, payload\_hex} и ссылок на DEK.
	\end{enumerate}
	
	Диаграмма сценария представлена на рисунке \ref{fig:load};
	
	\item сценарий \texttt{GET /data/\{name\}}
	
	\begin{enumerate}
	\item Поиск записи в \texttt{data\_records} по имени.
	\item Проверка \texttt{Authorization: Bearer} против \texttt{auth\_hash}.
	\item Получение \texttt{dek\_out}, извлечение ключа и валидация заголовка контейнера.
	\item Расшифрование \texttt{payload\_hex} и возврат открытого текста в API-ответе.
	\end{enumerate}
	
	Диаграмма сценария представлена на рисунке \ref{fig:get_data};
	
	\begin{figure}[p]
		\centering\includegraphics[height=0.95\textheight]{load.pdf}
    	\caption{Диаграмма сценария \texttt{load}}
    	\label{fig:load}
	\end{figure}
	
	\begin{figure}[p]
		\centering\includegraphics[height=0.95\textheight]{get_data.pdf}
	  \caption{Диаграмма сценария \texttt{GET /data/{name}}}
	  \label{fig:get_data}
	\end{figure}	
	
%	\begin{figure}[!htb]
%		\begin{minipage}{0.48\textwidth}
%		    \centering
%		    \includegraphics[height=0.7\textheight]{load.pdf}
%		    \caption{Диаграмма сценария \texttt{load}}
%		    \label{fig:load}
%		\end{minipage}
%		\hfill
%		\begin{minipage}{0.48\textwidth}
%		    \centering
%		    \includegraphics[height=0.7\textheight]{get_data.pdf}
%		    \caption{Диаграмма сценария \texttt{GET /data/{name}}}
%		    \label{fig:get_data}
%		\end{minipage}
%	\end{figure}
	
%	\FloatBarrier
	
	\end{itemize}
	
	\subsubsection{Обоснование выбора технологического стека}
	
	Выбор технологического стека для проекта сделан на основе двух критериев: 
	
	\begin{enumerate}
	\item обеспечение корректной и безопасной реализации криптографических операций.
	\item минимизация трудозатрат на разработку и верификацию в рамках дипломного цикла.
	\end{enumerate}
	
	Фактический стек реализации:
	
	\begin{enumerate}
	\item Python 3.11+ -- основной язык реализации CLI, API и сервисной логики.
	\item cryptography -- реализация AEAD-примитивов\footnote{См. раздел с определением AEAD-алгоритмов на с.~\pageref{def:AEAD}.} (AES-128-GCM, AES-256-GCM, ChaCha20-Poly1305) и криптоопераций на уровне production-библиотеки.
	\item FastAPI + uvicorn -- HTTP(S)-интерфейс и запуск API-сервера.
	\item Pydantic v2 + pydantic-settings -- строгая валидация входных данных, параметров CLI/API и конфигурации окружения.
	\item SQLAlchemy 2.x -- доступ к БД и реализация слоя хранения (\texttt{keys, data\_records}) через ORM.
	\item pytest -- модульное и интеграционное тестирование критичных сценариев.
	\end{enumerate}
	
	\subsubsection*{\textit{Обоснование выбора языка программирования}}
	
	Основным языком программирования для проекта выбран Python 3.13+, так как он обеспечивает высокий темп написания кода, имеет зрелую экосистему для backend и криптографии, а также удобен для тестирования и описания интерфейсов.
	
	Python характеризуется как хорошо читаемый, гибкий, платформонезависимый язык программирования, доступный для быстрого освоения как базового, так и профессионального уровня программирования. Он сочетает в себе возможности встроенных высокоуровневых структур данных и объектно-ориентированность с
	кратким и понятным синтаксисом. Благодаря своим качествам он
	подходит для решения широкого спектра задач, в том числе для автоматизации, анализа данных, машинного обучения, а также для
	программной реализации научных исследований и вычислений. Это
	универсальный и гибкий инструмент как для новичков, так и для
	опытных профессиональных разработчиков.
	Единственным недостатком Python как современного языка
	программирования стоит отметить сравнительно невысокую скорость выполнения программ, написанных на нем, что связано с их
	интерпретируемостью \cite{Sergeeva_python}.
	
	Сравнение с другими языками программирования показало, что для данного проекта Python имеет хороший баланс <<скорость реализации>>-<<читаемость кода>> и наличие библиотек под задачу конвертного шифрования. При этом ограничения Python, такие как интерпретируемость, зависимость от качества runtime-конфигурации, меньшая производительность в CPU-bound сценариях компенсируются архитектурными мерами: разделением KEK/DEK, хранением KEK вне БД, применением проверенных внешних библиотек и покрытием критичных ветвей тестами.
	
	
	
	На рисунке \ref{fig:languages} представлено сравнение языков программирования по сложности написания и поддержки с учетом специфики данного проекта.	
	
	\begin{figure}[H]
		\centering\includegraphics[width=\textwidth]{languages.pdf}
%		\centering\fitgraphics{languages.pdf}
		\caption{Сравнение языков программирования по сложности написания и поддержки}
		\label{fig:languages}
	\end{figure}	

	

	\FloatBarrier
	
	\subsubsection*{\textit{Обоснование выбора технологий}}
	
	Анализ выбранного технологического стека проводился с учётом требований к безопасности, производительности и удобству поддержки системы. Для каждой библиотеки были рассмотрены альтернативные решения, а также оценены их преимущества и ограничения в контексте данного проекта.
	
	Результаты сравнительного анализа представлены в таблице~\ref{tab:tech_stack_compare}.
	
	{\fontsize{12pt}{14pt}\selectfont
	\begin{tabularx}{\textwidth}{|p{3.8cm}|Z|Z|}
\caption{Сравнительный анализ выбранных технологий и альтернатив}
\label{tab:tech_stack_compare} \\

\hline
\multicolumn{1}{|c|}{Технология} &
\multicolumn{1}{c|}{Причины выбора} &
\multicolumn{1}{c|}{Альтернативы и ограничения} \\ \hline
\multicolumn{1}{|c|}{1} & 
\multicolumn{1}{c|}{2} &
\multicolumn{1}{c|}{3} \\ \hline
\endfirsthead

%\multicolumn{3}{r}{} \\[-0.3em]
\multicolumn{3}{r}{{\fontsize{14pt}{14pt}\selectfont Продолжение таблицы \thetable}} \\[0.5em]
\hline
%\multicolumn{1}{|c|}{Технология} &
%\multicolumn{1}{c|}{Причины выбора} &
%\multicolumn{1}{c|}{Альтернативы и ограничения} \\ \hline
\multicolumn{1}{|c|}{1} & 
\multicolumn{1}{c|}{2} &
\multicolumn{1}{c|}{3} \\ \hline
\endhead

uvicorn &
Лёгкий и производительный ASGI-сервер; простота запуска и интеграции с FastAPI; минимальная операционная сложность &
gunicorn -- более сложная конфигурация и ориентирован на production-оркестрацию; hypercorn -- менее распространён и имеет меньшую экосистему \\

\hline

cryptography &
Production-уровень реализации криптографических примитивов; поддержка AEAD-алгоритмов (AES-GCM, ChaCha20-Poly1305); использование OpenSSL; высокая надёжность и широкое применение в индустрии &
PyNaCl -- проще, но ограниченный набор алгоритмов; низкоуровневый OpenSSL API -- более гибкий, но существенно увеличивает риск ошибок реализации \\

\hline

FastAPI &
Нативная поддержка асинхронности (ASGI); высокая производительность; автоматическая генерация OpenAPI; тесная интеграция с Pydantic и строгая типизация &
Flask -- менее строгая типизация, требуется ручная валидация; Django REST Framework -- избыточен для лёгкого сервиса, выше порог сложности \\

\hline

Pydantic v2 + pydantic-settings &
Строгая валидация данных на основе аннотаций типов; высокая производительность (v2); единая модель данных для CLI, API и конфигурации; снижение количества ошибок &
marshmallow -- более многословен и менее удобен; dataclasses -- не предоставляют встроенной валидации и требуют ручной реализации \\

\hline

SQLAlchemy 2.x &
Зрелый ORM с гибкой архитектурой; поддержка как ORM, так и Core; контроль транзакций; независимость от фреймворка &
raw SQL -- больше контроля, но выше сложность и риск ошибок; Django ORM -- привязан к Django и избыточен для данного проекта \\

\hline

pytest &
Лаконичный синтаксис тестов; мощная система фикстур; удобство модульного и интеграционного тестирования; де-факто стандарт в Python &
unittest -- более громоздкий синтаксис и меньшая гибкость при организации тестов \\

\hline

\end{tabularx}	
	}
	Вывод: выбранный технологический стек выбран правильно с учетом специфики проекта и соответствует всем требованиям.
	
	\subsection{Модель угроз и допущений}
	
	Угроза безопасности информации -- это условия и факторы, создающие
	потенциальную или уже существующую опасность нарушения безопасности
	информации (ее конфиденциальности, целостности и доступности). Реализация
	угроз безопасности информации злоумышленником, как правило, возможна в
	информационных системах, автоматизированных системах управления,
	информационно-телекоммуникационных сетях, облачных инфраструктурах.
	
	Модель угроз (безопасности информации) -- это физическое,
	математическое и описательное представление свойств или характеристик угроз
	безопасности информации 
	
	Модель угроз -- это документ, назначение которого -
	определить угрозы, актуальные для конкретной системы \cite{Rayovskiy_Metody}.
	
	Модель угроз для данного проекта построена для архитектуры <<CLI/API + БД + криптографический модуль>> с учетом используемой схемы конвертного шифрования. Анализ выполняется в предположении, что основной ущерб связан с нарушением конфиденциальности хранимых данных и компрометацией ключевого материала.
	
	\subsubsection{Ключевые активы системы}
	
	В рамках рассматриваемой системы выделяется набор критически важных активов, компрометация которых может привести к раскрытию конфиденциальных данных или нарушению целостности системы. Основной упор делается на защите криптографических ключей, данных и механизмов аутентификации.
	
	\begin{enumerate}
	\item \texttt{VAULT\_MASTER\_KEK\_HEX} (master KEK), используемый для извлечения DEK.
	\item Обернутый DEK в таблице \texttt{keys}.
	\item Зашифрованные контейнеры и метаданные в \texttt{data\_records}.
	\item Заголовок авторизации (Authorization: Bearer) и его PBKDF2-представление.
	\item Исполняемая среда сервиса (процесс API/CLI, файловая система, переменные окружения).
	\end{enumerate}
	
	Перечисленные активы образуют основу модели защиты системы и определяют направления анализа потенциальных угроз.
		
	\subsubsection{Нарушители и их модель возможностей}
	
	Модель нарушителя строится с учётом различных уровней доступа к системе: от внешнего сетевого взаимодействия до компрометации инфраструктуры выполнения. Для каждого типа нарушителя предполагается различный набор возможностей и сценариев атак.
	
	\begin{enumerate}
	\item Внешний сетевой нарушитель: перехват/модификация трафика, попытки неавторизованного вызова API.
	\item Нарушитель с доступом к БД: чтение дампа таблиц, попытки восстановления открытого текста по содержимому БД.
	\item Нарушитель с компрометацией клиентского секрета: использование украденного Bearer-токена.
	\item Нарушитель с доступом к хосту выполнения: чтение переменных окружения и runtime-памяти процесса.
	\end{enumerate}
	
	Данная классификация позволяет системно рассмотреть угрозы, соответствующие различным уровням компрометации.
		
	\subsubsection{Границы доверия}
	
	Для корректного анализа угроз важно определить границы доверия, разделяющие компоненты системы по уровню контролируемости и потенциальной уязвимости.
	
	\begin{enumerate}
	\item Доверенная зона: процессы \texttt{vault/cli.py} и \texttt{vault/api.py}, выполняющие криптооперации и проверку авторизации.
	\item Условно доверенная зона: хранилище БД (\texttt{vault.db}), где данные считаются потенциально компрометируемыми.
	\item Недоверенная зона: сеть и внешние клиенты.
	\end{enumerate}
	
	Такое разделение позволяет применять различные механизмы защиты в зависимости от уровня доверия к компоненту.
		
	\subsubsection{Принятые допущения модели}
	
	При построении модели угроз были приняты следующие допущения, ограничивающие область анализа и определяющие базовые условия безопасности системы:
	
	\begin{enumerate}
	\item \texttt{VAULT\_MASTER\_KEK\_HEX} не хранится в БД и передается только через защищенное окружение выполнения.
	\item Хост сервиса администрируется корректно, а доступ к переменным окружения ограничен.
	\item TLS-сертификат и ключ сервера управляются безопасно и не скомпрометированы.
	\item Используются только поддерживаемые AEAD-алгоритмы и корректные размеры nonce/tag.
	\end{enumerate}
	
	Нарушение данных допущений существенно изменяет модель угроз и требует дополнительного анализа.
		
	\subsubsection{Основные сценарии угроз и контрмеры}
	
	На основе выделенных активов, модели нарушителя и границ доверия формируются ключевые сценарии угроз, а также соответствующие им меры защиты.
	
	\subsubsection*{\textit{Компрометация БД}}
	
	В данном сценарии предполагается получение злоумышленником доступа к содержимому базы данных.
	
	Контрмеры: DEK хранится только в обернутом виде; KEK отсутствует в БД; данные в \texttt{data\_records} сохраняются зашифрованными.
	
	\subsubsection*{\textit{Перехват/подмена сетевого трафика}}
	
	Сценарий предполагает возможность атаки типа Man-in-the-Middle и анализ сетевого взаимодействия.
	
	Контрмеры: HTTPS при запуске API; обязательная проверка Authorization; AEAD-проверка тега целостности при расшифровании.
	
	\subsubsection*{\textit{Компрометация Bearer-токена}}
	
	Предполагается утечка или перехват токена доступа клиента.
	
	Контрмеры: хранение не исходного токена, а PBKDF2-хеша; отказ в доступе при несовпадении; минимизация поверхности использования токена только в защищенном канале.
	
	\subsubsection*{\textit{Компрометация runtime-окружения}}
	
	Данный сценарий соответствует наиболее сильному нарушителю с доступом к системе выполнения.
	
	Контрмеры: вынос KEK в переменную окружения (не в БД), организационные меры hardening хоста, ограничение прав процесса. 
	
	\subsubsection{Диаграмма модели угроз}
	
	Финальная диаграмма модели угроз приложения показана на рисунке \ref{fig:threats}.
	
	\begin{figure}[!hbt]
		\centering\includegraphics[width=\textwidth]{threats.pdf}
		\caption{Модель угроз приложения}
		\label{fig:threats}
	\end{figure}
	
	\subsection{Схема взаимодействия с приложением}
	
	\subsubsection{Общая схема}
	
	Схема взаимодействия отражает два ключевых эксплуатационных сценария системы: загрузка и перешифрование данных через CLI, выдача данных через защищенный API по авторизации.
	
	В реализованном проекте взаимодействие построено вокруг принципа конвертного шифрования. Данные шифруются на уровне DEK (Data Encryption Key), а сами DEK не хранятся в открытом виде: перед записью в БД они оборачиваются master KEK (Key Encryption Key). Master KEK передается в процесс через переменную окружения \texttt{VAULT\_MASTER\_KEK\_HEX} и не сохраняется в таблицах приложения. Такое разделение уменьшает последствия компрометации базы данных: злоумышленник может получить зашифрованные записи и обернутый DEK, но без доступа к окружению выполнения не сможет корректно расшифровать полезную нагрузку.
	
	Участники процесса:
	
	\begin{enumerate}
	\item Администратор/разработчик, выполняющий команды управления и загрузки.
	\item CLI, реализующий операции \texttt{load}, работу с контейнером и запись в БД.
	\item Переменная окружения \texttt{VAULT\_MASTER\_KEK\_HEX}, из которой берется master KEK.
	\item Vault DB, хранящая обернутый DEK и зашифрованные записи.
	\item Vault API (HTTPS), предоставляющий endpoint чтения данных.
	\item Клиент приложения, запрашивающий данные через \texttt{GET /data/{name}}.
	\end{enumerate}
	
	Принципиальная особенность взаимодействия состоит в разделении ролей ключей: KEK никогда не сохраняется в БД, а DEK хранится только в обернутом виде. За счет этого даже при компрометации хранилища нарушитель не получает готовый ключ расшифрования без доступа к окружению выполнения.
	
	На рисунках \ref{fig:flow_load} и \ref{fig:flow_get} представлены диаграммы загрузки данных через CLI и получения данных через API соответственно.

	
	\begin{figure}[p]
		\centering\includegraphics[angle=90, height=0.95\textheight]{flow_load.pdf}
		\caption{Схема загрузки данных в приложение}
		\label{fig:flow_load}
	\end{figure}
	
	\begin{figure}[p]
		\centering\includegraphics[
		  width=0.9\textheight,
		  angle=90,
		  keepaspectratio
		]{flow_get.pdf}
		\caption{Схема получения данных из приложения}
		\label{fig:flow_get}
	\end{figure}	
		
	Интерпретация потока 1 (загрузка и перешифрование через CLI):
	
	\begin{enumerate}
	    \item Администратор инициирует команду \texttt{load} с указанием входного и выходного алгоритмов, а также имен DEK.
	    \item Входной контейнер может быть передан двумя способами: как JSON-файл со структурой header + data либо как нагрузка в шестнадцатеричном представлении вместе с отдельным \texttt{--header-json}. В обоих вариантах данные проходят валидацию через Pydantic-модели. В менее безопасном варианте данные могут передаваться через команду \texttt{store-plaintext} в открытом виде.
	    \item CLI извлекает из БД обернутые ключи \texttt{dek-in} и \texttt{dek-out}, после чего считывает master KEK из окружения.
	    \item Модуль \texttt{key\_manager}\listingfootnote{lst:vault/key_manager.py} выполняет извлечение обоих DEK. Если master KEK отсутствует, имеет неверную длину или не позволяет расшифровать обернутый DEK, операция прекращается без записи результата.
	    \item Модуль \texttt{crypto}\listingfootnote{lst:vault/crypto.py} проверяет заголовок входного контейнера: наличие полей \texttt{alg, nonce, tag}, соответствие ожидаемому алгоритму, размер nonce 12 байт и размер tag 16 байт.
	    \item После успешной проверки выполняется расшифрование входного контейнера алгоритмом \texttt{alg-in}, затем повторное шифрование открытого текста алгоритмом \texttt{alg-out}.
	    \item В БД сохраняется новый контейнер: \texttt{header\_json, payload\_hex}, имена использованных DEK, входной и выходной алгоритмы, а также \texttt{auth\_hash}. Открытый Bearer-токен не сохраняется.
	    \item Администратор получает структурированный JSON-ответ со статусом операции, выходным заголовком, размером открытого текста и сведениями о записи.
	\end{enumerate}
	
	Интерпретация потока 2 (получение данных через API):
	
	\begin{enumerate}
	    \item Клиент отправляет \texttt{GET /data/\{name\}} с заголовком Authorization: Bearer по HTTPS.
	    \item API ищет запись в \texttt{data\_records} по логическому имени. Если запись отсутствует, возвращается 404 Not Found.
	    \item Если заголовок Authorization отсутствует или не проходит проверку через \texttt{verify\_auth\_header}, выполнение прерывается с 401 Unauthorized.
	    \item После успешной авторизации API проверяет наличие заголовка контейнера и валидирует его относительно \texttt{alg\_out}, то есть относительно фактического алгоритма хранения данных после перешифрования.
	    \item API извлекает ключевую запись \texttt{dek\_out} из таблицы \texttt{keys}. Если ключ отсутствует или master KEK не передан в процесс, сервер возвращает ошибку уровня 500, поскольку нарушена внутренняя целостность конфигурации.
	    \item Master KEK нормализуется и проверяется, затем \texttt{key\_manager} выполняет извлечение целевого DEK.
	    \item Контейнер расшифровывается <<на лету>>. Открытый текст не записывается обратно в БД и существует только в памяти процесса во время обработки запроса.
	    \item Клиенту возвращается JSON-ответ с именем записи, алгоритмом, \texttt{plaintext\_hex}, опциональным \texttt{plaintext\_utf8} и числом байт. Если данные не являются валидной UTF-8 строкой, поле \texttt{plaintext\_utf8} возвращается как \texttt{null}, но шестнадцатеричное представление остается доступным.
	\end{enumerate}

	Обработка ошибок также является частью схемы взаимодействия. Ошибки пользовательского ввода на стороне CLI приводят к прекращению команды через \texttt{SystemExit} с диагностическим сообщением: например, при отсутствии файла, некорректном JSON, неверном шестнадцатеричном представлении или неизвестном DEK. Ошибки API преобразуются в HTTP-статусы: отсутствие записи соответствует 404, неверный Bearer -- 401, а повреждение внутренней конфигурации или контейнера -- 500. Такое разделение позволяет отличать ошибки клиента от нарушений целостности серверной конфигурации.
	
	Таким образом, схема взаимодействия демонстрирует сквозной защищенный конвейер обработки данных: от административной загрузки и перешифрования до контролируемой выдачи клиенту по API. Эта схема служит опорной для последующих разделов, где формализуются модель угроз, математическая модель конвертного шифрования и критерии тестирования.
	
	\subsubsection{Схема работы веб-сервера}
	
	Веб-сервер используется для сетевой выдачи данных через HTTPS API. В отличие от CLI-команд, которые выполняются как разовые административные операции, API-сервер должен запускаться как отдельный процесс и сохранять состояние своего запуска: адрес, порт, URL базы данных и идентификатор процесса. Поэтому для эксплуатации важны два симметричных сценария: безопасный старт сервера и корректная остановка ранее запущенного процесса.
	
	В проекте эти сценарии реализованы командами \texttt{web-server up} и \texttt{web-server down} в модуле \texttt{vault.cli}\listingfootnote{lst:vault/cli.py}. Команда запуска формирует процесс \texttt{uvicorn}, который поднимает приложение \texttt{vault.api:app} с TLS-ключом и сертификатом. Команда остановки использует служебный PID-файл \texttt{.vault\_web\_server.json}, чтобы найти ранее запущенный процесс и завершить его средствами операционной системы. Это решение не требует отдельного systemd-unit, Docker Compose или внешнего process manager, но при этом демонстрирует полный жизненный цикл API-сервиса.
	
	Для иллюстрации этих сценариев подготовлены две диаграммы последовательностей
	
	Первая диаграмма описывает процесс запуска сервера. Важно отметить, что запуск API является не просто вызовом uvicorn, а последовательностью проверок, которые защищают систему от конфликтов и некорректного состояния. CLI сначала валидирует параметры команды через \texttt{WebServerUpCommand}: проверяет существование TLS-ключа и сертификата, корректность host/port и допустимый диапазон порта. Затем выполняется проверка PID-файла \texttt{.vault\_web\_server.json}. Если файл существует и указанный в нем процесс еще активен, команда не запускает второй экземпляр сервера, а возвращает ошибку \texttt{already running}. Если PID-файл отсутствует, поврежден или указывает на уже завершенный процесс, он рассматривается как устаревший и удаляется.
	
	После проверки состояния CLI формирует команду запуска:

%	\vspace*{0.3cm}
	
	\needspace{2cm}
	
	\begin{lstlisting}[language=bash, caption={Строка запуска веб-сервера}, style=bashcommand]
python -m uvicorn vault.api:app --host <host> --port <port> --ssl-keyfile <key> --ssl-certfile <cert>\end{lstlisting}

	При запуске в окружение дочернего процесса передается \texttt{VAULT\_DB\_URL}, чтобы API использовал ту же БД, что и CLI-команды. Далее CLI ожидает короткий интервал и проверяет, не завершился ли процесс сразу после старта. Если процесс не смог запуститься, пользователь получает ошибку \texttt{Failed to start web server process.} Если запуск успешен, в \texttt{.vault\_web\_server.json} записываются \texttt{pid, host, port} и  \texttt{db\_url}, а CLI возвращает JSON-ответ со статусом ok и действием \texttt{web-server-up}.
	
	Вторая диаграмма описывает обратный сценарий -- остановку сервера. Ее следует рассматривать как эксплуатационную пару к схеме запуска: если \texttt{web-server up} создает фоновый процесс и PID-файл, то \texttt{web-server down} должен корректно обработать все возможные состояния этого файла. Сначала CLI проверяет наличие \texttt{.vault\_web\_server.json}. Если файла нет, команда не считает это ошибкой и возвращает результат \texttt{stopped=false} с причиной \texttt{not\_running}. Это удобно для повторного вызова команды остановки: она остается идемпотентной и не вызывает ошибок при создании сценариев автоматизации.
	
	Если PID-файл существует, CLI пытается извлечь из него идентификатор процесса. При поврежденном JSON или некорректном PID файл удаляется, а пользователю возвращается причина \texttt{invalid\_pid\_file}. Если PID валиден, CLI передает управление операционной системе: в Windows используется \texttt{taskkill /PID <pid> /T /F}, а в Linux/macOS применяется сигнал завершения через \texttt{os.kill(pid, 15)}. После этого CLI делает короткую паузу, проверяет, остался ли процесс активным, удаляет PID-файл и возвращает JSON-ответ \texttt{web-server-down} с признаком \texttt{stopped}.
	
	Диаграммы 1 и 2 представлены на рисунках \ref{fig:web-up} и \ref{fig:web-down} соответственно.
	
	\begin{figure}[!htb]
		\centering\includegraphics[angle=90, width=\textwidth]{web_up.pdf}
		\caption{Диаграмма запуска веб-сервера}
		\label{fig:web-up}
	\end{figure}


	\begin{figure}[!htb]
		\centering\includegraphics[angle=90, width=\textwidth]{web_down.pdf}
		\caption{Диаграмма остановки веб-сервера}
		\label{fig:web-down}
	\end{figure}
	
%	\FloatBarrier
	
	\subsubsection{Схема управления ключами}
	
	Для фиксации механизма управления ключами в приложении ниже приведены три отдельные диаграммы последовательностей. Разбиение делает явными разные фазы жизненного цикла ключевого материала: подготовку master KEK, перевод DEK в оберутое представление и последующее использование DEK в runtime.
	
	\subsubsection*{\textit{Подготовка master KEK}}
	
	Первая диаграмма описывает подготовку master KEK. В обоих вариантах, генерации и импорте, CLI либо создает значение ключа, либо нормализует уже существующий ключ в шестнадцатеричном формате, после чего возвращает его оператору без записи в таблицу \texttt{keys}. Завершающим шагом является перенос значения в \texttt{VAULT\_MASTER\_KEK\_HEX}, то есть в контур runtime-конфигурации, а не в контур долговременного хранения.
	
	Диаграмма представлена на рисунке \ref{fig:master_kek_prepare}.
	
	\begin{figure}[p]
		\centering\includegraphics[angle=90, height=0.95\textheight]{master_kek_prepare.pdf}
		\caption{Подготовка master KEK}
		\label{fig:master_kek_prepare}
	\end{figure}
	
	\subsubsection*{\textit{Создание или импорт DEK и его сохранение в обернутом виде}}
	
	Вторая диаграмма показывает цикл постановки DEK на хранение. Независимо от того, был ли ключ данных сгенерирован внутри приложения или импортирован из файла, перед сохранением он обязательно связывается с master KEK через операцию \texttt{wrap\_dek\_hex}. В результате в БД попадает не открытый DEK, а сериализованный обернутый пакет с \texttt{nonce}, шифртекстом и тегом аутентификации, что снижает последствия компрометации таблицы \texttt{keys}.
	
	Диаграмма представлена на рисунке \ref{fig:create_or_import_dek}.
	
	\begin{figure}[p]
		\centering\includegraphics[angle=90, height=0.95\textheight]{create_or_import_dek.pdf}
		\caption{Цикл постановки DEK на хранение}
		\label{fig:create_or_import_dek}
	\end{figure}	
	
	\subsubsection*{\textit{Использование DEK в runtime}}	
	
	Третья диаграмма отражает эксплуатационную фазу работы с ключом данных. Runtime-компонент CLI или API извлекает из БД обернутое представление DEK, считывает master KEK из окружения и выполняет \texttt{unwrap\_dek\_hex} через \texttt{key\_manager}. Открытый DEK появляется только в оперативной памяти процесса и используется немедленно для криптографической операции, после чего не сохраняется обратно в хранилище.
	
	Диаграмма представлена на рисунке \ref{fig:store_dek}.
	
	\begin{figure}[p]
		\centering\includegraphics[angle=90, height=0.95\textheight]{store_dek.pdf}
		\caption{Эксплуатационная фаза работы с ключом данных}
		\label{fig:store_dek}
	\end{figure}	

%	\FloatBarrier	
	
	\subsection{Обоснование выбора методики разработки}
		
	Чтобы успешно справляться с постоянно меняющимися потребностями и вызовами, с которыми сталкивается индустрия информационных технологий, разработчики и управляющие проектами активно развивают и применяют разнообразные методологии. История этих методологий богата и увлекательна, и она играет критическую
	роль в формировании современных подходов к разработке программного обеспечения и управлению IT-проектами \cite{Balanov_Agile}.
	
	
	\textbf{Методология TDD (Test-Driven-Development)} -- это подход к разработке
	ПО, который предполагает написание тестов перед написанием программного
	кода. Этот подход позволяет программистам убедиться, что их код работает
	корректно. 
	
	Основная идея TDD заключается в том, чтобы начинать с написания
	тестов для функциональности, которую вы планируете реализовать, а затем
	создавать код, который будет проходить эти тесты. Таким образом, тесты
	становятся не только инструментом для проверки работоспособности
	программного кода, но и своего рода спецификацией, определяющей
	требования к функционалу \cite{Shadrinsk_TDD}.
	
	Процесс TDD обычно следует циклу, известному как <<красный, зелёный,
	рефакторинг>>:
	
	\begin{enumerate}
	\item Красный (Red) -- Написание тестов, которые описывают ожидаемое
	поведение новой функциональности или исправление ошибок в существующем
	коде. На этом этапе тесты должны не проходить, что создаст <<красный>> сигнал \cite{Shadrinsk_TDD}.
	\item Зелёный (Green) -- Написание минимального кода, необходимого
	для прохождения написанных тестов. Цель -- сделать тесты проходящими и
	получить <<зелёный>> сигнал \cite{Shadrinsk_TDD}.
	\item Рефакторинг (Refactor) -- Улучшение кода без изменения его
	функциональности, чтобы сделать его более чистым, эффективным и
	поддерживаемым \cite{Shadrinsk_TDD}.
	\end{enumerate}
		
	Для реализации проекта была выбрана методика, близкая к TDD (Test-Driven Development). Для данного проекта такой подход особенно уместен, поскольку основные требования можно выразить как набор воспроизводимых инвариантов:
	
	\begin{enumerate}
	\item DEK не должен сохраняться в БД в открытом виде.
	\item Master KEK должен поступать из окружения и проходить проверку размера.
	\item Контейнер данных должен содержать корректные \texttt{alg, nonce, tag} и \texttt{ciphertext}.
	\item Расшифрование должно завершаться ошибкой при неверном ключе, алгоритме, теге или формате данных.
	\item API не должен возвращать открытый текст без корректного Authorization: Bearer.
	\item Хранимый Bearer-секрет должен быть представлен PBKDF2-хешем, а не исходной строкой.
	\item CLI-команды должны завершаться предсказуемо при ошибках входных параметров.
	\end{enumerate}
	
	На практике TDD в проекте использовался следующим образом: критичные сценарии сначала формулировались через тестируемое поведение, после чего реализовывались модули \texttt{crypto, key\_manager, auth, db, cli и api}. Такой подход позволил разделить разработку на короткие проверяемые итерации. Например, для криптографического слоя проверяется корректность шифрования и расшифрования контейнера, для авторизации -- хеширование и проверка Bearer-заголовка, для API -- успешная выдача данных и отказ при неверном токене.
	
	Альтернативой мог быть классический подход test-after development, при котором тесты пишутся после завершения реализации. Он проще на старте, однако хуже подходит для задач, где цена ошибки высока и требуется постоянно контролировать регрессии. Waterfall-подход также не является оптимальным: он предполагает жесткую последовательность этапов и хуже реагирует на уточнение контрактов данных, которое неизбежно возникает при разработке криптографического контейнера, CLI-команд и API. Итеративная Agile-модель близка по духу к выбранному подходу, но сама по себе не гарантирует тестируемости. Code-and-Fix обеспечивает максимальную скорость первого прототипа, однако для системы безопасности создает неприемлемый риск накопления скрытых дефектов.
	
	\subsection{Реализация ключевых модулей}
	
	Процесс разработки системы был организован вокруг постепенного наращивания функциональности: от минимального криптографического ядра к полноценному контуру CLI $\rightarrow$ БД $\rightarrow$ API. Такой порядок был выбран осознанно. В системе, работающей с ключами и зашифрованными данными, нельзя начинать с внешнего интерфейса и затем <<подключать безопасность>> как отдельный слой. Напротив, сначала необходимо зафиксировать корректное поведение базовых операций: проверку ключей, формат контейнера, шифрование, расшифрование, оборачивание DEK и отказ при некорректных данных.
	
	С учетом выбранной TDD-методики разработка велась короткими итерациями. Каждая итерация включала три шага: формулирование ожидаемого поведения, реализацию минимального кода для прохождения проверки и последующее уточнение структуры модуля. Это позволило не только быстрее обнаруживать ошибки, но и документировать саму логику системы через набор проверяемых сценариев. Например, требования к размеру nonce, формату Authorization, невозможности выдачи данных без Bearer и сохранению DEK только в обернутом виде были выражены не только в тексте, но и в тестируемом поведении.
	
	Разработка выполнялась в следующей последовательности:
	
	\begin{enumerate}
	\item Формирование криптографического ядра \texttt{vault/crypto.py}.
	\item Реализация управления ключами и конвертного шифрования в \texttt{vault/key\_manager.py}.
	\item Описание контрактов входных данных через Pydantic-схемы в \texttt{vault/schemas.py}\listingfootnote{lst:vault/schemas.py} и \texttt{vault/db\_schemas.py}\listingfootnote{lst:vault/db_schemas.py}.
	\item Реализация слоя хранения через SQLAlchemy в \texttt{vault/db.py}\listingfootnote{lst:vault/db.py}.
	\item Реализация авторизации и безопасного хранения Bearer-секрета в \texttt{vault/auth.py}\listingfootnote{lst:vault/auth.py}.
	\item Разработка CLI-команд в \texttt{vault/cli.py}.
	\item Разработка API-слоя в \texttt{vault/api.py}\listingfootnote{lst:vault/api.py}.
	\item Добавление интеграционных сценариев и сквозной проверки через pytest и \texttt{scripts/check\_all.py}\listingfootnote{lst:scripts/check_all.py}.
	\end{enumerate}
	
	Первым этапом стало создание криптографического модуля \texttt{vault/crypto.py}. Его задача состоит в том, чтобы изолировать непосредственные операции шифрования и расшифрования от остальной бизнес-логики. В модуле были зафиксированы поддерживаемые размеры ключей: 16 байт для AES-128-GCM, 32 байта для AES-256-GCM и 32 байта для ChaCha20-Poly1305. Это решение позволило не передавать ответственность за проверку длины ключа в CLI или API, а централизовать ее в одной функции \texttt{key\_bytes\_for\_algorithm}.
	
	Для контейнера данных был выбран компактный формат: отдельно хранится шифртекст в шестнадцатеричном виде и заголовок \texttt{header}, содержащий \texttt{alg, nonce} и \texttt{tag}. Такой формат удобен для записи в БД и для передачи между CLI/API. При этом заголовок не считается доверенным: перед расшифрованием он проходит проверку через \texttt{validate\_header}\calloutfullreffootnote{lst:func:validate_header}. Проверяются наличие обязательных полей, корректность шестнадцатеричных значений, совпадение алгоритма с ожидаемым, размер nonce 12 байт и размер AEAD-тега 16 байт. Если любое из условий нарушено, операция завершается ошибкой до попытки вернуть открытый текст.
	
	Следующим этапом была реализована работа с ключевым материалом в \texttt{vault/key\_manager.py}. Здесь было закреплено разделение ролей KEK и DEK. DEK используется для шифрования пользовательских данных, а master KEK используется для оборачивания DEK перед сохранением в БД. Функция \texttt{validate\_master\_kek\_hex}\calloutfullreffootnote{lst:func:validate_master_kek_hex} проверяет, что master KEK задан корректной шестнадцатеричной строкой и имеет размер 32 байта. Функция \texttt{wrap\_dek\_hex}\calloutfullreffootnote{lst:func:wrap_dek_hex} формирует JSON-структуру обернутого DEK с полями \texttt{v, wrap\_alg, nonce, wrapped} и \texttt{tag}; для оборачивания используется AES-256-GCM. Обратная операция \texttt{unwrap\_dek\_hex} восстанавливает исходный DEK только при наличии корректного master KEK.
		
	После криптографического слоя были описаны контракты данных. В \texttt{vault/schemas.py} описаны модели команд и входных контейнеров: \texttt{GenerateKeyCommand, AddKeyCommand, LoadCommand, WebServerUpCommand, HeaderModel, JsonContainerModel} и \texttt{FlatJsonContainerModel}. Это позволило перенести большую часть проверок с уровня ручных if-условий в декларативные Pydantic-модели. Например, \texttt{LoadCommand} проверяет существование входного файла, формат Bearer-заголовка, допустимость алгоритмов и зависимость между \texttt{--data-hex} и \texttt{--header-json}. Если данные передаются напрямую через \texttt{--data-hex}, заголовок становится обязательным, иначе система не сможет проверить алгоритм, nonce и tag.
	
	Для слоя хранения были созданы отдельные схемы \texttt{vault/db\_schemas.py}. Они валидируют данные перед записью в репозиторий и перед выдачей наружу. Такое разделение важно: схемы CLI описывают пользовательский ввод, а схемы БД описывают внутренние контракты хранения. Например, \texttt{DataRecordUpsertIn} проверяет \texttt{auth\_header}, алгоритмы, имена DEK, JSON-заголовок контейнера и шестнадцатеричное представление данных. 
	
	Слой хранения реализован в \texttt{vault/db.py} через SQLAlchemy. В нем определены две ORM-модели: \texttt{KeyRecord} и \texttt{DataRecord}. Таблица \texttt{keys} хранит логическое имя ключа, тип ключа, алгоритм и значение \texttt{key\_hex}. Для DEK это значение в целевом режиме содержит обернутый JSON. Таблица \texttt{data\_records} хранит имя записи, \texttt{auth\_hash}, входной и выходной алгоритмы, ссылки на \texttt{dek\_in} и \texttt{dek\_out}, \texttt{header\_json} и \texttt{payload\_hex}. В репозитории \texttt{VaultRepository} реализованы операции \texttt{upsert\_key}, \texttt{get\_key}, \texttt{upsert\_data\_record} и \texttt{get\_data\_record\_by\_name}\calloutfullreffootnote{lst:func:upsert_key}.
	
	Модуль авторизации \texttt{vault/auth.py} был выделен отдельно, чтобы не смешивать проверку Bearer с API-кодом или репозиторием. Функция \texttt{normalize\_bearer\_auth\_header}\calloutfullreffootnote{lst:func:normalize_bearer_auth_header} задает минимальный контракт заголовка: строка должна начинаться с Bearer и содержать непустой токен. Функция \texttt{hash\_auth\_header} сохраняет нормализованное значение в виде PBKDF2-хеша с алгоритмом \texttt{sha256}, солью 16 байт и 200000 итераций. Функция \texttt{verify\_auth\_header}\calloutfullreffootnote{lst:func:verify_auth_header} проверяет кандидатный заголовок через \texttt{hmac.compare\_digest}, что снижает риск утечки информации через различия во времени сравнения.
	
	После готовности сервисных модулей был реализован CLI-слой \texttt{vault/cli.py}. CLI стал основной административной точкой входа и получил команды \texttt{generate-dek}, \texttt{generate-kek}, \texttt{add-dek}, \texttt{add-kek}, \texttt{load}, \texttt{web-server up} и \texttt{web-server down}. Для каждой команды используется общий принцип: сначала аргументы разбираются через \texttt{argparse}, затем валидируются через Pydantic-модель, после чего вызываются сервисные функции и репозиторий. Результаты команд возвращаются в JSON-формате, что удобно для демонстрации, автоматизации и интеграционных проверок.
	
	Наиболее важной CLI-командой является \texttt{load}, поскольку она объединяет почти все компоненты системы. Команда получает входной контейнер, проверяет формат, извлекает \texttt{dek\_in} и \texttt{dek\_out} из БД, считывает master KEK из \texttt{VAULT\_MASTER\_KEK\_HEX}, выполняет извлечение обоих DEK, расшифровывает входной \texttt{payload} и повторно шифрует его целевым алгоритмом. После этого в \texttt{data\_records} сохраняется уже новый контейнер, а Bearer-заголовок передается в репозиторий только для хеширования и сохранения в поле \texttt{auth\_hash}.
	
	Завершающим прикладным слоем стал модуль \texttt{vault/api.py}. В нем реализована фабрика \texttt{create\_app}\calloutfullreffootnote{lst:func:create_app}, которая создает FastAPI-приложение, и endpoint \texttt{GET /data/{name}}. API выполняет последовательность проверок в безопасном порядке: сначала ищет запись, затем проверяет \texttt{Authorization}, затем валидирует заголовок контейнера, после этого получает \texttt{dek\_out}, выполняет извлечение через master KEK и только затем расшифровывает \texttt{payload}. Такой порядок принципиален: извлечение ключа и криптографическая операция не выполняются до успешной авторизации клиента. Процесс разработки производился как последовательное расширение ядра приложения. В начале были реализованы функции, отвечающие за криптографическую обработку данных. Затем появились контракты данных и слой хранения. После этого были добавлены команды CLI и сетевой API. Такой порядок позволил сохранить контролируемую сложность проекта: каждый новый слой опирался на уже проверенные нижележащие модули, а не дублировал их логику.	
	
	\newpage
	
	\section[Демонстрация результата работы]{ДЕМОНСТРАЦИЯ РЕЗУЛЬТАТА РАБОТЫ}
	
	Результат выполненной работы целесообразно показывать не только как перечень реализованных модулей, но и как воспроизводимый эксплуатационный сценарий. 
	
\subsubsection*{\textit{Генерация master KEK}}

Эта команда демонстрирует подготовку мастер-ключа, который используется не для шифрования пользовательских данных напрямую, а для оборачивания ключей данных. Важно подчеркнуть, что результат команды возвращается оператору в JSON-формате и не записывается в таблицу \texttt{keys}. Тем самым подтверждается один из ключевых инвариантов проекта: master KEK существует вне БД и должен передаваться в процесс только через защищенный канал конфигурации. Команда представлена на рисунке \ref{lst:generate_kek}.

\begin{lstlisting}[language=bash, style=bashcommand, caption={Создание KEK},label=lst:generate_kek]
PS> python main.py --db-url sqlite:///vault.db generate-kek --name master --alg AES-256-GCM
{
  "status": "ok",
  "action": "generate-kek",
  "algorithm": "AES-256-GCM",
  "key_hex": "592395cfaaaf4b42ec40f06336fe173e72b6a8137ec1abcc9b28b78467f80ee8",
  "stored_in_db": false,
  "note": "Store this KEK outside DB and set VAULT_MASTER_KEK_HEX for runtime operations."
}
\end{lstlisting}

\subsubsection*{\textit{Импорт ключей данных}}

Эти команды показывают постановку DEK на хранение. На вход CLI получает ключи в шестнадцатеричном формате из файлов, однако в БД попадают не исходные значения, а результат \texttt{wrap\_dek\_hex}, сформированный с использованием master KEK. Важно подчеркнуть, что логические имена <<\texttt{src}>> (рисунок \ref{lst:load_src_dek}) и <<\texttt{dst}>> (рисунок \ref{lst:load_dst_dek}) позволяют явно разделить входной и выходной контур шифрования при последующем перешифровании контейнера.

\needspace{1.5cm}

\begin{lstlisting}[style=bashcommand, caption={Создание src DEK}, label=lst:load_src_dek]
PS> python main.py --db-url sqlite:///vault.db add-dek --name src -f src_dek.hex
{
  "status": "ok",
  "action": "add-dek",
  "source_file": "src_dek.hex",
  "record": {
    "id": 2,
    "name": "src",
    "key_type": "dek",
    "algorithm": "IMPORTED",
    "key_bytes": 16,
    "created_at": "2026-03-01T05:35:34.545522",
    "updated_at": "2026-04-25T08:02:17.088711Z"
  },
  "dek_stored_as_wrapped": true
}
\end{lstlisting}

\begin{lstlisting}[style=bashcommand, caption={Создание dst DEK}, label=lst:load_dst_dek]
PS> python main.py --db-url sqlite:///vault.db add-dek --name dst -f dst_dek.hex
{
  "status": "ok",
  "action": "add-dek",
  "source_file": "dst_dek.hex",
  "record": {
    "id": 3,
    "name": "dst",
    "key_type": "dek",
    "algorithm": "IMPORTED",
    "key_bytes": 32,
    "created_at": "2026-03-01T05:35:35.108768",
    "updated_at": "2026-04-25T08:02:30.774061Z"
  },
  "dek_stored_as_wrapped": true
}\end{lstlisting}

\subsubsection*{\textit{Загрузка и перешифрование контейнера}}

Данная команда является ключевой демонстрацией результата работы системы, поскольку объединяет все основные подсистемы проекта. CLI считывает входной контейнер из \texttt{container.json}, проверяет структуру заголовка, извлекает из БД \texttt{src} и \texttt{dst} в обернутом виде, выполняет их извлечение через \texttt{VAULT\_MASTER\_KEK\_HEX}, расшифровывает данные исходным алгоритмом AES-128-GCM, затем повторно шифрует алгоритмом ChaCha20-Poly1305 и сохраняет новый контейнер в таблице \texttt{data\_records}. Одновременно \texttt{Authorization: Bearer 123} не сохраняется в открытом виде, а преобразуется в PBKDF2-хеш. Этот рисунок подтверждает достижение основной цели работы: централизованное управление ключами и контролируемое перешифрование данных. Команда представлена на рисунке \ref{lst:load_data}.

\begin{lstlisting}[style=bashcommand, caption={Загрузка данных из контейнера}, label=lst:load_data]
PS> python main.py --db-url sqlite:///vault.db load --name admin_data -f container.json --auth "Bearer 123" --alg-in AES-128-GCM --alg-out ChaCha20-Poly1305 --dek-in src --dek-out dst
{
  "status": "ok",
  "action": "load",
  "output_header": {
    "alg": "ChaCha20-Poly1305",
    "nonce": "a8b514fcc23093a22173349f",
    "tag": "f6aad059dd4b48ff91a16f8138273428"
  },
  "plaintext_bytes": 11,
  "record": {
    "id": 3,
    "name": "admin_data",
    "alg_in": "AES-128-GCM",
    "alg_out": "ChaCha20-Poly1305",
    "dek_in": "src",
    "dek_out": "dst",
    "auth_protected": true,
    "payload_bytes": 11,
    "has_header": true,
    "created_at": "2026-04-25T08:02:57.816990Z",
    "updated_at": "2026-04-25T08:02:57.816995Z"
  }
}
\end{lstlisting}

\subsubsection*{\textit{Загрузка и перешифрование открытых данных в произвольной форме}}

Программа также поддерживает загрузку данных произвольной формы из файла. Для этого используется команда \texttt{store-plaintext}. Дальнейшее взаимодействие с такими данными ничем не отличается от основного метода взаимодействия с программой. Команда представлена на рисунке \ref{lst:load_data_plain}.

\begin{lstlisting}[style=bashcommand, caption={Загрузка данных произвольного формата}, label=lst:load_data_plain]
PS> python main.py --db-url sqlite:///vault.db store-plaintext --name admin_data2 -f .\sample.txt --auth "Bearer 123" --alg ChaCha20-Poly1305 --dek dst
{
  "status": "ok",
  "action": "store-plaintext",
  "output_header": {
    "alg": "ChaCha20-Poly1305",
    "nonce": "9263a2ebececb77e4d785752",
    "tag": "c960a5f4b351e7ba19cd2984f43e2d93"
  },
  "plaintext_bytes": 7,
  "record": {
    "id": 4,
    "name": "admin_data2",
    "alg_in": "ChaCha20-Poly1305",
    "alg_out": "ChaCha20-Poly1305",
    "dek_in": "dst",
    "dek_out": "dst",
    "auth_protected": true,
    "payload_bytes": 7,
    "has_header": true,
    "created_at": "2026-04-25T09:08:58.294473Z",
    "updated_at": "2026-04-25T09:08:58.294477Z"
  }
}
\end{lstlisting}

\subsubsection*{\textit{Запуск защищенного API-сервера}}

Здесь демонстрируется перевод подготовленных данных в режим сетевой выдачи. Команда проверяет наличие TLS-ключа и сертификата, запускает \texttt{uvicorn} в фоновом процессе и сохраняет служебную информацию о процессе в \texttt{.vault\_web\_server.json}. Административная загрузка данных и их последующая выдача намеренно разделены: сначала запись формируется через CLI, затем уже существующая запись читается клиентом через минимальный HTTPS endpoint. Команда представлена на рисунке \ref{lst:web_up}.



\begin{lstlisting}[style=bashcommand, caption={Запуск веб-сервера}, label=lst:web_up]
PS> python main.py --db-url sqlite:///vault.db web-server up --key server.key --cert server.crt --host 127.0.0.1 --port 9000
{
  "status": "ok",
  "action": "web-server-up",
  "pid": 16468,
  "host": "127.0.0.1",
  "port": 9000
}
\end{lstlisting}

\subsubsection*{\textit{Проверка доступности сервиса}}

Этот вспомогательный запрос показывает, что веб-сервер успешно поднят и готов обслуживать клиентские обращения. Он не затрагивает криптографические операции и не требует авторизации, поэтому его удобно использовать как промежуточную проверку корректности запуска перед основной демонстрацией выдачи данных. Команда представлена на рисунке \ref{lst:web_check}.

\begin{lstlisting}[style=bashcommand, caption={Проверка состояния веб-сервера}, label=lst:web_check]
PS> curl.exe -k https://127.0.0.1:9000/health
{"status":"ok"}
\end{lstlisting}

\subsubsection*{\textit{Выдача данных по авторизованному запросу}}

Этот запрос является итоговой демонстрацией прикладного результата работы. API находит запись \texttt{admin\_data} в таблице \texttt{data\_records}, проверяет Bearer-заголовок относительно сохраненного \texttt{auth\_hash}, извлекает \texttt{dek\_out} из таблицы \texttt{keys}, считывает master KEK из окружения, выполняет извлечение ключа данных и расшифровывает контейнер непосредственно перед возвратом ответа. В ответ клиент получает JSON со значениями \texttt{name, algorithm, plaintext\_hex}, опциональным \texttt{plaintext\_utf8} и размером полезной нагрузки в байтах. Таким образом подтверждается, что сервис не только хранит зашифрованные данные, но и выдает их только при корректной авторизации и наличии согласованной ключевой конфигурации. Команда представлена на рисунке \ref{lst:api_get}.

\begin{lstlisting}[style=bashcommand, caption={Выдача данных через API}, label=lst:api_get]
PS> curl.exe -k -H "Authorization: Bearer 123" https://127.0.0.1:9000/data/admin_data
{"name":"admin_data","algorithm":"ChaCha20-Poly1305","plaintext_hex":"68656c6c6f2d7661756c74","plaintext_utf8":"hello-vault","bytes":11}
\end{lstlisting}

\subsubsection*{\textit{Остановка API-сервера}}

Финальная команда завершает сценарий и показывает, что жизненный цикл веб-сервера также управляется средствами самого приложения. CLI считывает PID-файл, завершает фоновый процесс и удаляет служебное состояние. Команда показана на рисунке \ref{lst:web_down}.

\begin{lstlisting}[style=bashcommand, caption={Остановка веб-сервера}, label=lst:web_down]
PS> python main.py web-server down
{
  "status": "ok",
  "action": "web-server-down",
  "stopped": true,
  "pid": 16468
}
\end{lstlisting}

\subsubsection*{\textit{Создание контейнера}}

Для глубокой интеграции с процессом разработки для создания контейнеров был выбран подход создания программного API для шифрования контейнеров вместо традиционного создания CLI команд. API реализован в виде функции \texttt{encrypt\_container} из модуля \texttt{vault.crypto}. Пример скрипта, который создает зашифрованный контейнер представлен на рисунке \ref{lst:container_api} (Приложение Б, с. \pageref{lst:container_api}).

Сценарии, в которых такой подход оправдан:

\begin{itemize}
\item внутренний сервис хранения чувствительных данных. Приложение сохраняет не открытый текст, а контейнер, а потом получает данные обратно только по авторизованному запросу через \texttt{GET /data/{name}};
\item backend для нескольких микросервисов. Один сервис отвечает за создание и хранение контейнеров, а другие не работают с ключами напрямую, а обращаются к Vault API за расшифрованными данными;
\item защищенное хранение конфигурации или технических секретов. Например, токены интеграций, ключи внешних API, параметры доступа к стендам. Их можно хранить в виде контейнеров и выдавать только авторизованным клиентам;
\item локальные и тестовые контуры разработки. Когда команде нужен не тяжелый KMS, а компактный сервис, который дает единый процесс: положить данные, зашифровать, хранить, выдать по токену.
\end{itemize}

\clearpage
	
	\newpage
	
	\genericsection{Заключение}
	
	В рамках дипломной работы была рассмотрена задача разработки защищенного сервиса управления криптографическими ключами и зашифрованными данными. Актуальность выбранной темы определяется тем, что современные программные системы используют большое количество чувствительных данных: ключи шифрования, токены доступа, технические учетные данные и конфиденциальные пользовательские сведения. При отсутствии централизованного и воспроизводимого механизма работы с такими данными возрастают риски их утечки, некорректного хранения, случайной публикации в репозиториях и несанкционированного доступа.
	
	Цель работы заключалась в проектировании и реализации прикладного сервиса Vault, который обеспечивает безопасную обработку зашифрованных данных, управление ключевым материалом и авторизованную выдачу данных через API. Поставленная цель была достигнута: разработан программный комплекс, включающий консольный интерфейс администратора, криптографический модуль, модуль управления ключами, слой хранения на базе SQLAlchemy, FastAPI-приложение и набор автоматизированных проверок.
	
	В ходе анализа предметной области были рассмотрены существующие подходы и решения для управления секретами и ключами: HashiCorp Vault, облачные KMS/Secrets Manager-сервисы, PAM/Secrets-платформы, self-hosted решения и GitOps-инструменты шифрования конфигураций. По результатам анализа был обоснован выбор собственного специализированного решения для учебно-прикладного проекта. Такой подход позволил сосредоточиться на ключевых механизмах безопасности: конвертное шифрование, разделении KEK и DEK, контроле формата криптографического контейнера, авторизации запросов и воспроизводимости операций через CLI.
	
	В проектной части была разработана модульно-слоистая архитектура системы. Входными точками стали \texttt{vault/cli.py} и \texttt{vault/api.py}, сервисный слой представлен модулями \texttt{vault/crypto.py}, \texttt{vault/key\_manager.py} и \texttt{vault/auth.py}, а доступ к данным изолирован в \texttt{VaultRepository}. Такое разделение позволило снизить связанность компонентов: CLI и API используют общую криптографическую и валидационную логику, а работа с БД не дублируется в прикладных сценариях.
	
	Ключевым архитектурным решением стала модель конвертного шифрования. В реализованной системе данные шифруются с использованием DEK, а сами DEK сохраняются в базе данных только в обернутом виде. Master KEK не записывается в БД и передается в runtime через переменную окружения \texttt{VAULT\_MASTER\_KEK\_HEX}. Благодаря этому компрометация базы данных сама по себе не дает злоумышленнику готового материала для расшифрования данных. В БД остаются зашифрованные контейнеры, обернутый DEK, метаданные алгоритмов и PBKDF2-хеши авторизационных заголовков.
	
	Криптографическая часть реализована на основе AEAD-алгоритмов AES-128-GCM, AES-256-GCM и ChaCha20-Poly1305. Использование AEAD позволило объединить конфиденциальность и контроль целостности: при повреждении шифротекста, неверном ключе, неправильном nonce, несоответствии алгоритма или некорректном теге расшифрование завершается ошибкой. Для контейнера данных были зафиксированы обязательные поля \texttt{alg, nonce, tag} и \texttt{payload} в шестнадцатеричном представлении, а проверка этих полей вынесена в отдельный слой валидации.
	
	Важным результатом работы стала реализация двух основных эксплуатационных сценариев. Первый сценарий связан с загрузкой и перешифрованием данных через CLI: администратор передает входной контейнер, указывает входной и выходной алгоритмы, имена DEK и Bearer-секрет, после чего система выполняет извлечение ключей, расшифрование, повторное шифрование и сохранение нового контейнера. Второй сценарий связан с выдачей данных через \texttt{GET /data/{name}}: API проверяет Authorization, извлекает запись, выполняет извлечение нужного DEK и расшифровывает данные только после успешной авторизации.
	
	Также внимание было уделено безопасному хранению авторизационных данных. Bearer-заголовок не сохраняется в БД в открытом виде. Вместо этого используется PBKDF2-хеширование с солью и 200000 итераций. Проверка заголовка выполняется через \texttt{hmac.compare\_digest}, что снижает риск утечки информации через различия во времени сравнения.
	
	Для реализации проекта была выбрана методика разработки, близкая к TDD. Это решение оказалось обоснованным, поскольку безопасность системы зависит от набора строгих инвариантов: DEK не должен храниться открыто, API не должен выдавать открытый текст без валидного Bearer, контейнер должен проходить проверку структуры, а ошибки криптографических операций не должны приводить к раскрытию данных. Автоматизированные тесты позволили закрепить эти инварианты и контролировать регрессии при изменении модулей.
	
	Проверка решения выполнялась с помощью unit-тестов, интеграционных тестов и сквозного сценария \texttt{scripts/check\_all.py}. Проверялись криптографические roundtrip-сценарии, валидация схем, работа CLI, API-доступ, отказ при неверной авторизации, миграционная совместимость БД и полный поток загрузки данных. Дополнительно был выполнен ручной тест, включающий создание ключей, формирование входного контейнера, перешифрование через CLI и чтение результата через API. По результатам проверки автоматизированные тесты и интеграционный сценарий завершились успешно.

	\newpage
	\genericsection{Библиографический список}
	\vspace*{-0.3cm}

	\begin{refcontext}[sorting=customsort]
	\nocite{*}
	\printbibliography[heading=none]
	\end{refcontext}
		
	\label{bib:end}

	\afterpage{\setlength{\footskip}{3cm}}
	\afterpage{\setlength{\skip\footins}{\baselineskip}}
%	\afterpage{\listingsgeometry}

	\newpage

	\newgeometry{
	left=3cm,
	right=1cm,
	top=2cm,
	bottom=2cm,
	includehead
	}

	\newappendix

\phantomsection
\addcontentsline{toc}{section}{Приложение \currentappendixletter}

%\setlength{\footskip}{2cm}

\thispagestyle{plain}
\pagestyle{appendixcontstyle}

\vspace*{-\headheight}
\vspace*{-\headsep}
\vspace*{-0.8cm}

\centering{ПРИЛОЖЕНИЕ \currentappendixletter}

\vspace*{0.3cm}

\lstset{style=modern} 

\lstinputlisting[caption=main.py,label=lst:main.py,style=modern]{main.py}
\lstinputlisting[caption=pytest.ini,label=lst:pytest.ini]{pytest.ini}
\lstinputlisting[caption=requirements.txt,label=lst:requirements.txt]{requirements.txt}
\lstinputlisting[caption=vault/\_\_init\_\_.py,label=lst:vault/__init__.py]{vault/\_\_init\_\_.py}
\lstinputlisting[caption=vault/api.py,label=lst:vault/api.py]{vault/api.py}
\lstinputlisting[caption=vault/auth.py,label=lst:vault/auth.py]{vault/auth.py}	
\lstinputlisting[caption=vault/cli.py,label=lst:vault/cli.py]{vault/cli.py}	
\lstinputlisting[caption=vault/crypto.py,label=lst:vault/crypto.py]{vault/crypto.py}	
\lstinputlisting[caption=vault/db.py,label=lst:vault/db.py]{vault/db.py}	
\lstinputlisting[caption=vault/db\_schemas.py,label=lst:vault/db_schemas.py]{vault/db\_schemas.py}
\lstinputlisting[caption=vault/key\_manager.py,label=lst:vault/key_manager.py]{vault/key\_manager.py}	
\lstinputlisting[caption=vault/schemas.py,label=lst:vault/schemas.py]{vault/schemas.py}		
\lstinputlisting[caption=vault/settings.py,label=lst:vault/settings.py]{vault/settings.py}		
\lstinputlisting[caption=tests/conftest.py,label=lst:tests/conftest.py]{tests/conftest.py}		
\lstinputlisting[caption=tests/test\_api.py,label=lst:tests/test_api.py]{tests/test\_api.py}		
\lstinputlisting[caption=tests/test\_auth.py,label=lst:tests/test_auth.py]{tests/test\_auth.py}		
\lstinputlisting[caption=tests/test\_cli.py,label=lst:tests/test_cli.py]{tests/test\_cli.py}		
\lstinputlisting[caption=tests/test\_crypto.py,label=lst:tests/test_crypto.py]{tests/test\_crypto.py}		
\lstinputlisting[caption=tests/test\_db\_migration.py,label=lst:tests/test_db_migration.py]{tests/test\_db\_migration.py}		
\lstinputlisting[caption=tests/test\_load\_flow.py,label=lst:tests/test_load_flow.py]{tests/test\_load\_flow.py}
\lstinputlisting[caption=tests/test\_schemas.py,label=lst:tests/test_schemas.py]{tests/test\_schemas.py}
\lstinputlisting[caption=scripts/check\_all.py,label=lst:scripts/check_all.py]{scripts/check\_all.py}
	\newpage
	\newappendix

\phantomsection
\addcontentsline{toc}{section}{Приложение \currentappendixletter}

%\setlength{\footskip}{2cm}

\thispagestyle{plain}
\pagestyle{appendixcontstyle}

\centering{ПРИЛОЖЕНИЕ \currentappendixletter}

\vspace*{0.3cm}

\lstset{style=modern} 

\lstinputlisting[language=python, style=modern, caption={Пример использования программного API для создания зашифрованного контейнера}, label=lst:container_api]{scripts/create_container_demo.py}

\end{document}
